\documentclass[11pt]{article}

\usepackage[a4paper,margin=1in]{geometry}
\usepackage{amsmath,amssymb,amsthm,mathtools}
\usepackage{bm}
\usepackage{mathrsfs}
\usepackage{hyperref}
\usepackage{enumitem}
\usepackage{booktabs}
\usepackage{microtype}

\hypersetup{
  colorlinks=true,
  linkcolor=blue,
  citecolor=blue,
  urlcolor=blue
}

\newtheorem{theorem}{Theorem}
\newtheorem{proposition}[theorem]{Proposition}
\newtheorem{lemma}[theorem]{Lemma}
\newtheorem{corollary}[theorem]{Corollary}
\theoremstyle{definition}
\newtheorem{definition}[theorem]{Definition}
\newtheorem{assumption}[theorem]{Assumption}
\newtheorem{remark}[theorem]{Remark}

\newcommand{\R}{\mathbb{R}}
\newcommand{\E}{\mathbb{E}}
\newcommand{\Prob}{\mathbb{P}}
\newcommand{\supp}{\mathrm{supp}}
\newcommand{\norm}[1]{\left\lVert #1 \right\rVert}
\newcommand{\abs}[1]{\left\lvert #1 \right\rvert}
\newcommand{\ip}[2]{\left\langle #1, #2 \right\rangle}
\newcommand{\mc}[1]{\mathcal{#1}}
\newcommand{\cE}{\mathcal{E}}
\newcommand{\cB}{\mathcal{B}}

\title{Finite-Horizon Mean-Field Training and Linear Mode Connectivity\\for a Single-Gaussian Natural-Softplus Head}
\author{Working draft}
\date{\today}

\begin{document}
\maketitle

\begin{abstract}
We study a two-layer network that outputs a single Gaussian predictive
distribution through the raw head
\[
  z_\Theta(x)=\bigl(u_\Theta(x),s_\Theta(x)\bigr)
\]
and the natural map
\[
  \eta_\Theta(x)=\bigl(u_\Theta(x),-\psi_{\rm sp}(s_\Theta(x))\bigr).
\]
The Gaussian loss is convex in natural parameters, but parameter interpolation
is not the same as natural-output interpolation because of the softplus scale
link.

Our training-side contribution is finite-horizon: for noiseless SGD, two
width-$N$ runs started from the same initialization law track the same
deterministic mean-field path on every fixed interval $[0,T]$. This is a
shared-path statement, not a long-time convergence theorem to a stationary
law. Our interpolation-side contribution is a deterministic aligned barrier
theorem for the practical single-Gaussian head, together with sharper raw-space
and exact conditional-risk refinements collected in Appendix~E. Combining the
common-path statement with Wasserstein alignment yields a finite-horizon LMC
theorem for the practical single-Gaussian head. Appendices~A--C contain the
main proofs, Appendix~D records rate refinements, and Appendix~E records the
sharpened deterministic barrier route.
\end{abstract}

\section{Introduction}

Linear mode connectivity (LMC) refers to the empirical phenomenon that two
independently trained networks can often be connected by a low-loss path after
an appropriate neuron permutation. This behavior was documented empirically,
for example, by Garipov et al.~\cite{Garipov2018} and
Frankle et al.~\cite{Frankle2020}. In the two-layer mean-field setting,
Ferbach et al.~\cite{Ferbach2024} gave the first quantitative explanation for
convex Lipschitz losses. Their argument has a clean two-stage structure:
\begin{enumerate}[label=(\roman*)]
  \item a training-side statement showing that two wide runs started from the
  same law remain close to the same deterministic mean-field trajectory;
  \item an interpolation-side statement showing that, after optimal transport
  alignment of the hidden neurons, the loss barrier along the parameter
  interpolation path is small.
\end{enumerate}

This paper develops that program for heteroscedastic single-Gaussian
regression. The head is practice faithful but no longer scalar. The network
produces a two-coordinate raw output
\[
  z_\Theta(x)=\bigl(u_\Theta(x),s_\Theta(x)\bigr)\in\R^2,
\]
and the actual Gaussian natural parameters are
\[
  \eta_\Theta(x)=\Gamma(z_\Theta(x))
  =
  \bigl(u_\Theta(x),-\psi_{\rm sp}(s_\Theta(x))\bigr),
\]
where $\psi_{\rm sp}(s)=\lambda_{\min}+\log(1+e^s)$ keeps the second natural
coordinate negative. Here $u_\Theta$ is the first natural coordinate itself,
not the predictive mean; the latter is
\[
  \mu_\Theta(x)=\frac{u_\Theta(x)}{2\psi_{\rm sp}(s_\Theta(x))}.
\]
This is the model one would train directly in raw parameters.

The single-Gaussian case is the clean place to begin. The head is already
genuinely vector valued, yet the output family remains elementary enough to
track exactly. The Gaussian negative log-likelihood is exactly convex in
natural coordinates, but the practical softplus implementation breaks affine
interpolation in natural space. The interpolation problem therefore survives in
a simple but still nontrivial form.

The contributions are threefold.
\begin{enumerate}[label=(C\arabic*)]
  \item We prove a finite-horizon mean-field approximation and a shared-law
  common-path theorem for noiseless SGD with the practical natural-softplus
  single-Gaussian head. The training-side novelty is that this vector-valued,
  nonlinearly reparameterized Gaussian head still fits the finite-horizon SGD
  mean-field framework after a model-specific drift, confinement, and
  regularity analysis.
  \item We prove a deterministic aligned barrier theorem in natural-parameter
  space. The barrier is controlled by two explicit defects: vector-head
  transfer and softplus curvature, and Appendix~E shows how this route can be
  sharpened through raw-space and exact conditional-risk reductions.
  \item Combining the shared-path statement with optimal-transport alignment
  yields an LMC theorem for the practical $K=1$ distributional model.
\end{enumerate}

The training-side statement is deliberately finite horizon. We do not prove
long-time convergence to a stationary law. What we prove is that, on each fixed
interval $[0,T]$, two wide runs started from the same law follow the same
deterministic distributional path. This is exactly the training-side input
needed for the LMC bridge. For readability, the main text states the theorem
stack first. Appendices~A--C contain the model-specific verification and full
proofs, Appendix~D collects rate refinements, and Appendix~E records sharper
deterministic barrier refinements and quantitative instantiations.

\section{Setup}

We separate the barrier-side assumptions from the additional second-order
feature regularity needed on the training side.

\begin{assumption}[Barrier-side regularity]
\label{ass:barrier-k1}
Let $\mc{X}\subseteq\R^d$ be an input space and let $(X,Y)$ be a data pair. We
assume $|Y|\le Y_\ast$ almost surely. For notational simplicity we take the
hidden parameter to be $\xi\in\R^d$; this is the standard case without an
explicit bias, and bias terms can be folded into $d$ by the usual augmentation.
Fix a feature map $\varphi:\R^d\times\mc{X}\to\R$ satisfying
\[
  K_\varphi:=\sup_{\xi,x}|\varphi(\xi,x)|<\infty,
  \qquad
  K_{\nabla}:=\sup_{\xi,x}\norm{\nabla_\xi \varphi(\xi,x)}<\infty.
\]
\end{assumption}

\begin{assumption}[Second-order feature regularity]
\label{ass:feature-second-k1}
Under Assumption~\ref{ass:barrier-k1}, assume furthermore that
$\nabla_\xi\varphi(\cdot,x)$ is globally Lipschitz in $\xi$, uniformly in $x$,
i.e.
\[
  \norm{\nabla_\xi\varphi(\xi,x)-\nabla_\xi\varphi(\bar\xi,x)}
  \le
  L_{\nabla}\norm{\xi-\bar\xi}
  \qquad
  \forall \xi,\bar\xi\in\R^d,\ \forall x\in\mc{X},
\]
for some finite constant $L_{\nabla}$.
\end{assumption}

\begin{remark}
Both assumptions are stated for the composite feature map $\varphi$, not
directly for an activation. In the common single-index case
\[
  \varphi(\xi,x)=\sigma(\langle\xi,x\rangle),
  \qquad
  \sup_{x\in\mc{X}}\norm{x}\le R_X<\infty,
\]
a sufficient condition for Assumption~\ref{ass:barrier-k1} is that $\sigma$ is
bounded and $\sigma'$ is bounded. Assumption~\ref{ass:feature-second-k1} is
further implied by $\sigma'$ being globally Lipschitz on $\R$ (equivalently, if
$\sigma\in C^2$, by $\sigma''$ being bounded). In that case
\[
  K_\varphi\le \|\sigma\|_\infty,
  \qquad
  K_{\nabla}\le R_X\|\sigma'\|_\infty,
  \qquad
  L_{\nabla}\le R_X^2\mathrm{Lip}(\sigma').
\]
Thus the barrier theorem covers bounded activations with bounded derivative on a
bounded input domain, while the training-side mean-field theorem additionally
asks for a globally Lipschitz derivative. Smooth bounded activations such as
$\tanh$, logistic sigmoid, and arctangent satisfy both assumptions, whereas
ReLU, leaky-ReLU, GELU, and softplus are excluded as hidden activations.
\end{remark}

\begin{remark}
The interpolation-side barrier theorems
(Theorems~\ref{thm:main-barrier-k1} and \ref{thm:deterministic-lmc-k1}) use
only Assumption~\ref{ass:barrier-k1}. The training-side mean-field theorems
(Theorems~\ref{thm:main-shared-path-k1}, \ref{thm:sgd-mean-field-k1}, and
\ref{thm:main-sgd-lmc-k1}) additionally use
Assumption~\ref{ass:feature-second-k1}; this extra assumption enters through
the stochastic-drift Lipschitz estimate in
Proposition~\ref{prop:sample-drift-regularity-k1}. The refined $O(B_N)$ rate in
Appendix~D also uses Assumption~\ref{ass:feature-second-k1}. The sharpened
deterministic barrier appendix (Appendix~E) uses only the barrier-side
notation, together with additional endpoint or pathwise quantities when stated.
\end{remark}

\subsection{The practical single-Gaussian head}

We consider width-$N$ two-layer networks with mean-field scaling
\begin{equation}
  z_\Theta(x)
  =
  \bigl(u_\Theta(x),s_\Theta(x)\bigr)
  :=
  \frac1N\sum_{i=1}^N (a_i,c_i)\,\varphi(\xi_i,x)
  \in \R^2,
  \label{eq:raw_head}
\end{equation}
where
\[
  \Theta=((a_i,c_i,\xi_i))_{i=1}^N\in(\R\times\R\times\R^d)^N.
\]
The natural output is
\begin{equation}
  \eta_\Theta(x)
  :=
  \Gamma(z_\Theta(x))
  =
  \bigl(u_\Theta(x),-\psi_{\rm sp}(s_\Theta(x))\bigr),
  \label{eq:natural_head}
\end{equation}
with
\[
  \psi_{\rm sp}(s):=\lambda_{\min}+\log(1+e^s),
  \qquad
  \lambda_{\min}>0.
\]
Equivalently, the induced Gaussian has conditional mean and variance
\[
  \mu_\Theta(x)=\frac{u_\Theta(x)}{2\psi_{\rm sp}(s_\Theta(x))},
  \qquad
  \sigma_\Theta^2(x)=\frac{1}{2\psi_{\rm sp}(s_\Theta(x))}.
\]
Thus $u_\Theta$ is the first natural coordinate, not the predictive mean
itself.

In the corresponding mean-field description, a measure
$\rho\in\mc{P}(\R^2\times\R^d)$ induces
\begin{align}
  \eta_{1,\rho}(x)
  &:=
  \int a\,\varphi(\xi,x)\,\rho(d\theta),
  \\
  s_\rho(x)
  &:=
  \int c\,\varphi(\xi,x)\,\rho(d\theta),
  \\
  \eta_{2,\rho}(x)
  &:=
  -\psi_{\rm sp}(s_\rho(x)),
\end{align}
where $\theta=(a,c,\xi)$.

\begin{proposition}[Automatic negativity of the second natural coordinate]
\label{prop:softplus-domain-k1}
For every $\rho$ and every $x\in\mc{X}$,
\[
  \eta_{2,\rho}(x)\le -\lambda_{\min}<0.
\]
The same pointwise bound holds for every finite-width realization
\eqref{eq:natural_head}.
\end{proposition}

\begin{proof}
By definition, $\psi_{\rm sp}(s)\ge \lambda_{\min}$ for all $s\in\R$.
\end{proof}

\subsection{Gaussian natural loss and natural-strip bounds}

For Gaussian natural parameters
$\eta=(\eta_1,\eta_2)\in \R\times(-\infty,0)$, define
\[
  A_{\rm G}(\eta_1,\eta_2)
  =
  -\frac{\eta_1^2}{4\eta_2}
  +
  \frac12\log\!\Bigl(\frac{-\pi}{\eta_2}\Bigr),
\]
and
\begin{equation}
  \ell_\eta(\eta;y)
  :=
  A_{\rm G}(\eta)-\eta_1 y-\eta_2 y^2.
  \label{eq:gaussian_nll_eta}
\end{equation}
The population risk is
\begin{equation}
  \cE(\Theta)
  :=
  \E\big[\ell_\eta(\eta_\Theta(X);Y)\big].
  \label{eq:population_risk}
\end{equation}

For later use define
\begin{align}
  q_1(\eta;y)
  &:=
  \partial_{\eta_1}\ell_\eta(\eta;y)
  =
  -\frac{\eta_1}{2\eta_2}-y,
  \label{eq:q1-k1}
  \\
  q_2(\eta;y)
  &:=
  \partial_{\eta_2}\ell_\eta(\eta;y)
  =
  \frac{\eta_1^2}{4\eta_2^2}
  -
  \frac{1}{2\eta_2}
  -
  y^2.
  \label{eq:q2-k1}
\end{align}

\begin{proposition}[Convexity of the Gaussian natural loss]
\label{prop:gaussian-natural-convexity}
For every fixed $y\in\R$, the map
\[
  \eta\mapsto \ell_\eta(\eta;y)=A_{\rm G}(\eta)-\eta_1 y-\eta_2 y^2
\]
is convex on the half-space $\{\eta_2<0\}$.
\end{proposition}

\begin{proof}
The Gaussian log-partition $A_{\rm G}$ is convex on $\{\eta_2<0\}$ because it
is the log-partition function of a regular minimal exponential family; the
linear term preserves convexity. See, for example,
Wainwright and Jordan~\cite{WainwrightJordan2008}.
\end{proof}

\begin{proposition}[Gaussian strip smoothness]
\label{prop:gaussian-smoothness-k1}
Fix $M>0$ and $\varepsilon>0$. On the strip
\[
  D_{M,\varepsilon}
  :=
  \{(\eta_1,\eta_2)\in\R^2: |\eta_1|\le M,\ \eta_2\le -\varepsilon\},
\]
the Gaussian log-partition function has globally bounded Hessian:
\[
  \sup_{(\eta_1,\eta_2)\in D_{M,\varepsilon}}
  \norm{\nabla^2 A_{\rm G}(\eta_1,\eta_2)}_{\mathrm{op}}
  \le
  L_{\rm G}(M,\varepsilon),
\]
where one may take
\[
  L_{\rm G}(M,\varepsilon)
  :=
  \frac{1}{2\varepsilon}
  +
  \frac{M}{2\varepsilon^2}
  +
  \frac{M^2}{2\varepsilon^3}
  +
  \frac{1}{2\varepsilon^2}.
\]
\end{proposition}

\begin{proof}
The Hessian entries are
\begin{align*}
  \partial_{11}^2 A_{\rm G}(\eta)
  &=
  -\frac{1}{2\eta_2},
  \\
  \partial_{12}^2 A_{\rm G}(\eta)
  &=
  \frac{\eta_1}{2\eta_2^2},
  \\
  \partial_{22}^2 A_{\rm G}(\eta)
  &=
  -\frac{\eta_1^2}{2\eta_2^3}+\frac{1}{2\eta_2^2}.
\end{align*}
Since $\eta_2\le -\varepsilon$ and $|\eta_1|\le M$, each entry is bounded in
absolute value by the corresponding term in $L_{\rm G}(M,\varepsilon)$.
\end{proof}

\begin{proposition}[Convexity and Lipschitz control on a natural strip]
\label{prop:convex-lipschitz-strip-k1}
Fix $M>0$ and $\lambda_{\min}>0$, and define
\[
  D_{M,\lambda_{\min}}
  :=
  \{(\eta_1,\eta_2)\in\R^2: |\eta_1|\le M,\ \eta_2\le -\lambda_{\min}\}.
\]
Then for every $|y|\le Y_\ast$:
\begin{enumerate}[label=(\roman*)]
  \item $\ell_\eta(\cdot;y)$ is convex on $D_{M,\lambda_{\min}}$;
  \item $\ell_\eta(\cdot;y)$ is
  $C_\eta(M,\lambda_{\min},Y_\ast)$-Lipschitz in the $\ell_\infty$ norm on
  $D_{M,\lambda_{\min}}$, where
  \begin{equation}
    C_\eta(M,\lambda_{\min},Y_\ast)
    :=
    \frac{M}{2\lambda_{\min}} + Y_\ast
    + \frac{M^2}{4\lambda_{\min}^2}
    + \frac{1}{2\lambda_{\min}}
    + Y_\ast^2.
    \label{eq:Ceta-k1}
  \end{equation}
\end{enumerate}
\end{proposition}

\begin{proof}
Convexity follows from Proposition~\ref{prop:gaussian-natural-convexity}.
For the Lipschitz bound, write
\[
  q_1(\eta;y)
  :=
  \partial_{\eta_1}\ell_\eta(\eta;y)
  =
  -\frac{\eta_1}{2\eta_2}-y
\]
and
\[
  q_2(\eta;y)
  :=
  \partial_{\eta_2}\ell_\eta(\eta;y)
  =
  \frac{\eta_1^2}{4\eta_2^2}
  -
  \frac{1}{2\eta_2}
  -
  y^2.
\]
On $D_{M,\lambda_{\min}}$ and $|y|\le Y_\ast$,
\[
  |q_1(\eta;y)|
  \le
  \frac{M}{2\lambda_{\min}}+Y_\ast,
  \qquad
  |q_2(\eta;y)|
  \le
  \frac{M^2}{4\lambda_{\min}^2}
  +
  \frac{1}{2\lambda_{\min}}
  +
  Y_\ast^2.
\]
The mean-value theorem in the $\ell_\infty$ norm proves the claim.
\end{proof}

\subsection{Aligned interpolation and barrier}

Let $\Theta_A$ and $\Theta_B$ be two trained width-$N$ networks. Fix a neuron
permutation $\varpi\in\mathfrak S_N$ and absorb it into network $B$, so we write
\[
  \Theta_A=((a_i^A,c_i^A,\xi_i^A))_{i=1}^N,
  \qquad
  \Theta_B=((a_i^B,c_i^B,\xi_i^B))_{i=1}^N.
\]
The parameter interpolation path is
\[
  \Theta_t
  :=
  \bigl((1-t)a_i^A+ta_i^B,\ (1-t)c_i^A+tc_i^B,\ (1-t)\xi_i^A+t\xi_i^B\bigr)_{i=1}^N.
\]
We also define
\[
  z_A:=z_{\Theta_A},
  \qquad
  z_B:=z_{\Theta_B},
  \qquad
  z_t:=z_{\Theta_t},
\]
the raw-output interpolation
\[
  \hat z_t:=(1-t)z_A+tz_B,
\]
the natural-output interpolation
\[
  \hat\eta_t:=(1-t)\eta_{\Theta_A}+t\eta_{\Theta_B},
\]
and the lifted raw interpolation
\[
  \tilde\eta_t:=\Gamma(\hat z_t).
\]

\begin{definition}[Aligned barrier]
For a fixed neuron permutation $\varpi$, define
\[
  \cB(\Theta_A,\Theta_B;\varpi)
  :=
  \sup_{t\in[0,1]}
  \Bigl\{
    \cE(\Theta_t)
    -
    (1-t)\cE(\Theta_A)
    -
    t\cE(\Theta_B)
  \Bigr\}.
\]
\end{definition}

For later use define the hidden mismatch
\begin{equation}
  B_N
  :=
  \frac1N\sum_{i=1}^N \norm{\xi_i^A-\xi_i^B}^2,
  \label{eq:hidden_mismatch_k1}
\end{equation}
the endpoint raw-scale gap
\begin{equation}
  \Delta_{s,N}
  :=
  \E\big[(s_{\Theta_A}(X)-s_{\Theta_B}(X))^2\big],
  \label{eq:raw_scale_gap}
\end{equation}
and the endpoint head moment
\begin{equation}
  M_V^\infty
  :=
  \max\!\Biggl\{
    \Bigl(\frac1N\sum_{i=1}^N (a_i^A)^2\Bigr)^{1/2},
    \Bigl(\frac1N\sum_{i=1}^N (c_i^A)^2\Bigr)^{1/2},
    \Bigl(\frac1N\sum_{i=1}^N (a_i^B)^2\Bigr)^{1/2},
    \Bigl(\frac1N\sum_{i=1}^N (c_i^B)^2\Bigr)^{1/2}
  \Biggr\}.
  \label{eq:head_moment_k1}
\end{equation}

\section{Main Results}

This section states the three theorem-level contributions. The main text keeps
only the theorem stack and short proof roadmaps. The model-specific drift
analysis, confinement estimates, interpolation lemmas, and Wasserstein-alignment
arguments are deferred to Appendices~A--E.

The interpolation-side results below use Assumption~\ref{ass:barrier-k1}. The
training-side results additionally use
Assumption~\ref{ass:feature-second-k1} together with the following
finite-horizon SGD regime.

\begin{assumption}[Finite-horizon SGD regime]
\label{ass:softplus-sgd-k1}
Fix a horizon $T<\infty$. The step sizes satisfy
\[
  s_k=\varepsilon\,\xi(k\varepsilon),
\]
where $\xi:[0,\infty)\to(0,\infty)$ is bounded and Lipschitz, and the data
stream
\[
  Z_{k+1}=(X_{k+1},Y_{k+1})
\]
is i.i.d.\ with law $P$. The initialization law $\rho_0$ on
$\R^2\times\R^d$ has compact support.
\end{assumption}

\subsection{Training-Side Shared Path for SGD}

\begin{theorem}[Training-side shared mean-field path for SGD]
\label{thm:main-shared-path-k1}
Assume Assumptions~\ref{ass:barrier-k1}, \ref{ass:feature-second-k1}, and
\ref{ass:softplus-sgd-k1}. Then there exist a deterministic measure-valued path
$(\rho_t)_{t\in[0,T]}$ and an error term
$\delta_{N,\varepsilon}(T)\to 0$ as $N\to\infty$ and $\varepsilon\downarrow 0$
such that the following holds. For any two independent width-$N$ noiseless SGD
runs for the practical natural-softplus Gaussian head, with empirical laws
$\hat\rho_{A,k}^N$ and $\hat\rho_{B,k}^N$, both initialized i.i.d.\ from the
same law $\rho_0$, one has with high probability
\[
  \sup_{k\le T/\varepsilon}
  W_1\!\bigl(\hat\rho_{A,k}^N,\rho_{k\varepsilon}\bigr)
  \le
  \delta_{N,\varepsilon}(T),
  \qquad
  \sup_{k\le T/\varepsilon}
  W_1\!\bigl(\hat\rho_{B,k}^N,\rho_{k\varepsilon}\bigr)
  \le
  \delta_{N,\varepsilon}(T).
\]
Consequently,
\[
  \sup_{k\le T/\varepsilon}
  W_1\!\bigl(\hat\rho_{A,k}^N,\hat\rho_{B,k}^N\bigr)
  \le
  2\,\delta_{N,\varepsilon}(T)
\]
with high probability, and the left-hand side converges to zero in probability
as $N\to\infty$ and $\varepsilon\downarrow 0$.
\end{theorem}

\begin{proof}[Proof sketch]
Appendix~A derives the raw stochastic drift of the practical natural-softplus
Gaussian head, proves finite-horizon confinement and uniform drift regularity,
and checks the bounded-support hypotheses required by the
Mei--Misiakiewicz--Montanari/Ferbach finite-horizon SGD machinery. The stated
shared-path theorem is then exactly the combination of
Theorem~\ref{thm:sgd-mean-field-k1} and
Corollary~\ref{cor:sgd-common-path-k1}.
\end{proof}

\subsection{Interpolation-Side Deterministic Barrier}

\begin{theorem}[Interpolation-side deterministic barrier theorem]
\label{thm:main-barrier-k1}
Assume Assumption~\ref{ass:barrier-k1}. For every fixed neuron permutation
$\varpi$,
\begin{equation}
  \cB(\Theta_A,\Theta_B;\varpi)
  \le
  C_\eta(K_\varphi M_V^\infty,\lambda_{\min},Y_\ast)
  \left[
    \frac12 K_{\nabla}M_V^\infty\sqrt{B_N}
    +
    \frac1{32}\Delta_{s,N}
  \right].
  \label{eq:main-deterministic-barrier-k1}
\end{equation}
The two terms on the right are exactly the two model-specific interpolation
defects: the vector-head transfer term and the softplus Jensen term. In
particular, if $B_N\to 0$ and $\Delta_{s,N}\to 0$, then
\[
  \cB(\Theta_A,\Theta_B;\varpi)\to 0.
\]
\end{theorem}

\begin{proof}[Proof sketch]
Appendix~B proves the two explicit ingredients behind
\eqref{eq:main-deterministic-barrier-k1}: a vector-head transfer bound and a
softplus Jensen bound. Exact convexity of the Gaussian loss in natural
parameters then turns those two defects into the deterministic barrier bound.
The full argument is recorded in Theorem~\ref{thm:deterministic-lmc-k1}.
\end{proof}

\begin{remark}
Appendix~E records sharper deterministic refinements of
Theorem~\ref{thm:main-barrier-k1}. Those results replace the final
natural-strip Lipschitz reduction by a raw-space reduction, exact
sigmoid-projected transfer functionals, and exact conditional-risk moduli.
They are not needed for the main finite-horizon common-path to LMC chain, but
they materially tighten theorem-consistent instantiated certificates on the
matched softplus run.
\end{remark}

\subsection{Main Finite-Horizon SGD-LMC Theorem}

\begin{theorem}[Finite-horizon SGD linear mode connectivity]
\label{thm:main-sgd-lmc-k1}
Assume Assumptions~\ref{ass:barrier-k1}, \ref{ass:feature-second-k1}, and
\ref{ass:softplus-sgd-k1}. Let $\hat\rho_{A,k}^N$ and $\hat\rho_{B,k}^N$ be the
empirical laws of two independent width-$N$ SGD runs of the practical
natural-softplus Gaussian head, with parameter vectors $\Theta_{A,k}$ and
$\Theta_{B,k}$, both initialized i.i.d.\ from the same law $\rho_0$ and
trained over the same horizon $[0,T]$.
Let
$A_T^{\rm sgd},C_T^{\rm sgd},R_T^{\rm sgd}$ be the deterministic confinement
envelopes furnished in Appendix~A, and write
\[
  M_{V,T}^{\rm sgd}:=\max\{A_T^{\rm sgd},C_T^{\rm sgd}\}.
\]
Then there exist permutations $\varpi_{N,k}$ such that, with high probability and
uniformly for all $k\le T/\varepsilon$,
\begin{align}
  \cB(\Theta_{A,k},\Theta_{B,k};\varpi_{N,k})
  &\le
  C_\eta(K_\varphi M_{V,T}^{\rm sgd},\lambda_{\min},Y_\ast)
  \Biggl[
    \frac12 K_{\nabla}M_{V,T}^{\rm sgd}\sqrt{2R_T^{\rm sgd}}\,
    \bigl(2\delta_{N,\varepsilon}(T)\bigr)^{1/2}
    \nonumber\\
  &\qquad\qquad
    +
    \frac{(K_\varphi+C_T^{\rm sgd}K_{\nabla})^2}{8}\,
    \delta_{N,\varepsilon}(T)^2
  \Biggr].
  \label{eq:main-sgd-lmc-k1}
\end{align}
In particular, for every fixed finite horizon, the aligned interpolation
barrier vanishes in probability as $N\to\infty$ and $\varepsilon\downarrow 0$.
\end{theorem}

\begin{proof}[Proof sketch]
Theorem~\ref{thm:main-shared-path-k1} provides a common deterministic
mean-field path for the two SGD runs. Appendix~C turns the resulting
Wasserstein control into a neuron permutation and converts the endpoint mismatch
into the alignment quantities $B_N$ and $\Delta_{s,N}$. Applying
Theorem~\ref{thm:main-barrier-k1} then yields the stated LMC bound; the
detailed reduction is recorded in Theorem~\ref{thm:lmc-from-common-path-k1}.
\end{proof}

\section{Conclusion}

The paper has three theorem-level contributions: a finite-horizon shared
mean-field path for SGD, a deterministic natural-parameter barrier theorem, and
the resulting finite-horizon SGD-LMC theorem for the practical single-Gaussian
natural-softplus head. This is the cleanest way to expose the logic of the
result: training-side common path, interpolation-side barrier, then LMC.

Appendices~A--C contain the model-specific verification and full proofs behind
that theorem stack. Appendix~D records rate refinements, and Appendix~E records
sharper deterministic barrier refinements and quantitative instantiations.

\appendix

\section{Training-Side Verification and Proofs}

This appendix proves Theorem~\ref{thm:main-shared-path-k1}. We record the
model-specific raw drift, confinement, and regularity estimates needed to place
the practical natural-softplus Gaussian head inside the bounded-support
finite-horizon SGD mean-field framework of
Mei--Misiakiewicz--Montanari~\cite{MeiMisiakiewiczMontanari2019} and
Ferbach et al.~\cite{Ferbach2024}.

\subsection{Deterministic drift and finite-horizon confinement}

\begin{proposition}[Raw drift formula]
\label{prop:softplus-raw-drift-k1}
For $\theta=(a,c,\xi)\in\R^2\times\R^d$ and a probability measure $\rho$,
define
\[
  \eta_\rho(x)
  :=
  \bigl(\eta_{1,\rho}(x),\eta_{2,\rho}(x)\bigr),
  \qquad
  h_\rho(x,y)
  :=
  \psi_{\rm sp}'(s_\rho(x))\,q_2(\eta_\rho(x);y).
\]
Then the mean-field-scaled population gradient flow of the risk
\eqref{eq:population_risk} is the interacting particle system
\begin{equation}
  \dot\theta_i^N(t)
  =
  b\bigl(\theta_i^N(t),\rho_t^N\bigr),
  \qquad
  \rho_t^N:=\frac1N\sum_{i=1}^N \delta_{\theta_i^N(t)},
  \label{eq:particle-system-k1}
\end{equation}
with drift $b=(b_a,b_c,b_\xi)$ given by
\begin{align}
  b_a(\theta,\rho)
  &:=
  -
  \E\big[q_1(\eta_\rho(X);Y)\,\varphi(\xi,X)\big],
  \label{eq:drift-a-k1}
  \\
  b_c(\theta,\rho)
  &:=
  \E\big[h_\rho(X,Y)\,\varphi(\xi,X)\big],
  \label{eq:drift-c-k1}
  \\
  b_\xi(\theta,\rho)
  &:=
  -
  \E\Big[
    \big(a\,q_1(\eta_\rho(X);Y)-c\,h_\rho(X,Y)\big)\nabla_\xi\varphi(\xi,X)
  \Big].
  \label{eq:drift-xi-k1}
\end{align}
\end{proposition}

\begin{proof}
Let
\[
  \mc{L}_N(\theta_1,\dots,\theta_N)
  :=
  \E\big[\ell_\eta(\eta_{\rho_N}(X);Y)\big],
  \qquad
  \rho_N=\frac1N\sum_{j=1}^N \delta_{\theta_j}.
\]
Since
\[
  \partial_{\theta_i}\eta_{1,\rho_N}(x)
  =
  \frac1N\bigl(\varphi(\xi_i,x),0,a_i\nabla_\xi\varphi(\xi_i,x)\bigr)
\]
and
\[
  \partial_{\theta_i}\eta_{2,\rho_N}(x)
  =
  -\frac1N\psi_{\rm sp}'(s_{\rho_N}(x))
  \bigl(0,\varphi(\xi_i,x),c_i\nabla_\xi\varphi(\xi_i,x)\bigr),
\]
the chain rule yields
\[
  \nabla_{\theta_i}\mc{L}_N
  =
  \frac1N
  \begin{pmatrix}
    \E\big[q_1(\eta_{\rho_N}(X);Y)\varphi(\xi_i,X)\big]
    \\
    -\E\big[h_{\rho_N}(X,Y)\varphi(\xi_i,X)\big]
    \\
    \E\big[(a_iq_1(\eta_{\rho_N}(X);Y)-c_i h_{\rho_N}(X,Y))
    \nabla_\xi\varphi(\xi_i,X)\big]
  \end{pmatrix}.
\]
Multiplying by the mean-field factor $N$ in continuous-time gradient flow gives
\eqref{eq:particle-system-k1}--\eqref{eq:drift-xi-k1}.
\end{proof}

The associated deterministic mean-field equation is
\begin{equation}
  \partial_t \rho_t
  +
  \nabla_\theta\cdot\bigl(\rho_t\,b(\theta,\rho_t)\bigr)
  =
  0,
  \qquad
  \rho_{t=0}=\rho_0.
  \label{eq:mf-pde-k1}
\end{equation}

\begin{proposition}[Finite-horizon confinement]
\label{prop:softplus-confinement-k1}
Fix a horizon $T<\infty$ and let
\[
  A_0:=\sup_{(a,c,\xi)\in\supp(\rho_0)}|a|,
  \qquad
  C_0:=\sup_{(a,c,\xi)\in\supp(\rho_0)}|c|,
  \qquad
  R_0:=\sup_{(a,c,\xi)\in\supp(\rho_0)}\norm{\xi}.
\]
Define the common initialization box
\[
  B_0
  :=
  \{(a,c,\xi): |a|\le A_0,\ |c|\le C_0,\ \norm{\xi}\le R_0\}.
\]
Define envelopes on $[0,T]$ by
\begin{align}
  \dot A_t
  &=
  K_\varphi\left(\frac{K_\varphi A_t}{2\lambda_{\min}}+Y_\ast\right),
  \qquad
  A_{t=0}=A_0,
  \label{eq:A-envelope-k1}
  \\
  \dot C_t
  &=
  K_\varphi\left(
    \frac{K_\varphi^2 A_t^2}{4\lambda_{\min}^2}
    +
    \frac{1}{2\lambda_{\min}}
    +
    Y_\ast^2
  \right),
  \qquad
  C_{t=0}=C_0,
  \label{eq:C-envelope-k1}
  \\
  \dot R_t
  &=
  K_{\nabla}\left[
    A_t\left(\frac{K_\varphi A_t}{2\lambda_{\min}}+Y_\ast\right)
    +
    C_t\left(
      \frac{K_\varphi^2 A_t^2}{4\lambda_{\min}^2}
      +
      \frac{1}{2\lambda_{\min}}
      +
      Y_\ast^2
    \right)
  \right],
  \qquad
  R_{t=0}=R_0.
  \label{eq:R-envelope-k1}
\end{align}
Then $A_t$, $C_t$, and $R_t$ remain finite on $[0,T]$. Define for each
$t\in[0,T]$
\[
  B_t
  :=
  \{(a,c,\xi): |a|\le A_t,\ |c|\le C_t,\ \norm{\xi}\le R_t\},
  \qquad
  M_t:=K_\varphi A_t.
\]
Every characteristic of any solution path of \eqref{eq:mf-pde-k1} whose
initial law is supported on $B_0$, as well as every particle path of
\eqref{eq:particle-system-k1} whose initial particles lie in $B_0$, satisfies
\[
  |a(t)|\le A_t,
  \qquad
  |c(t)|\le C_t,
  \qquad
  \norm{\xi(t)}\le R_t
  \qquad
  \forall t\in[0,T].
\]
In particular, for every such mean-field solution path,
\[
  \supp(\rho_t)\subseteq B_t\subseteq B_T
  \qquad
  \forall t\in[0,T],
\]
and
\[
  |\eta_{1,\rho_t}(x)|\le M_t\le M_T
  \qquad
  \forall x\in\mc{X},\ \forall t\in[0,T].
\]
\end{proposition}

\begin{proof}
Because $\eta_{2,\rho}(x)\le -\lambda_{\min}$ and
$|\eta_{1,\rho}(x)|\le K_\varphi A_t$, \eqref{eq:q1-k1}--\eqref{eq:q2-k1}
give
\[
  |q_1(\eta_\rho(x);y)|
  \le
  \frac{K_\varphi A_t}{2\lambda_{\min}}+Y_\ast
\]
and
\[
  |q_2(\eta_\rho(x);y)|
  \le
  \frac{K_\varphi^2 A_t^2}{4\lambda_{\min}^2}
  +
  \frac{1}{2\lambda_{\min}}
  +
  Y_\ast^2.
\]
Since $0<\psi'_{\rm sp}<1$, the drift formulas imply
\[
  |\dot a(t)|\le \dot A_t,
  \qquad
  |\dot c(t)|\le \dot C_t,
  \qquad
  \norm{\dot\xi(t)}\le \dot R_t.
\]
Scalar comparison yields the claim.
\end{proof}

\subsection{SGD mean-field approximation and a shared-law common path}

Throughout this subsection, assume Assumptions~\ref{ass:barrier-k1} and
\ref{ass:softplus-sgd-k1}. We invoke
Assumption~\ref{ass:feature-second-k1} only when second-order feature
regularity is needed.

\begin{definition}[One-sample stochastic drift]
\label{def:sample-drift-k1}
For $z=(x,y)\in \mc{X}\times\R$, define
\[
  \hat b(\theta,\rho;z)
  :=
  \bigl(\hat b_a(\theta,\rho;z),\hat b_c(\theta,\rho;z),\hat b_\xi(\theta,\rho;z)\bigr)
\]
by
\begin{align}
  \hat b_a(\theta,\rho;z)
  &:=
  -q_1(\eta_\rho(x);y)\,\varphi(\xi,x),
  \label{eq:sample-drift-a-k1}
  \\
  \hat b_c(\theta,\rho;z)
  &:=
  \psi'_{\rm sp}(s_\rho(x))\,q_2(\eta_\rho(x);y)\,\varphi(\xi,x),
  \label{eq:sample-drift-c-k1}
  \\
  \hat b_\xi(\theta,\rho;z)
  &:=
  -
  \bigl(
    a\,q_1(\eta_\rho(x);y)
    -
    c\,\psi'_{\rm sp}(s_\rho(x))q_2(\eta_\rho(x);y)
  \bigr)\nabla_\xi\varphi(\xi,x).
  \label{eq:sample-drift-xi-k1}
\end{align}
\end{definition}

\begin{proposition}[Unbiased raw stochastic update]
\label{prop:unbiased-sample-drift-k1}
Under Assumption~\ref{ass:softplus-sgd-k1},
\[
  b(\theta,\rho)=\E_Z[\hat b(\theta,\rho;Z)],
\]
where $Z=(X,Y)\sim P$. Consequently the noiseless SGD iteration can be written
as
\begin{equation}
  \theta_i^{k+1}
  =
  \theta_i^k
  +
  s_k\,\hat b(\theta_i^k,\hat\rho_k^N;Z_{k+1}),
  \qquad
  \hat\rho_k^N:=\frac1N\sum_{i=1}^N \delta_{\theta_i^k}.
  \label{eq:sgd-update-k1}
\end{equation}
\end{proposition}

\begin{proof}
Take the expectation of
\eqref{eq:sample-drift-a-k1}--\eqref{eq:sample-drift-xi-k1}
with respect to $Z=(X,Y)\sim P$ and compare with
\eqref{eq:drift-a-k1}--\eqref{eq:drift-xi-k1}.
\end{proof}

\begin{proposition}[Time-inhomogeneous finite-horizon confinement]
\label{prop:sgd-confinement-k1}
Assume Assumptions~\ref{ass:barrier-k1} and \ref{ass:softplus-sgd-k1}, and let
\[
  \xi_{\max}:=\sup_{t\in[0,T]}\xi(t)<\infty.
\]
Define $A_t^{\rm sgd}$, $C_t^{\rm sgd}$, and $R_t^{\rm sgd}$ exactly as in
\eqref{eq:A-envelope-k1}--\eqref{eq:R-envelope-k1}, but with the right-hand
sides multiplied by $\xi_{\max}$. Then every characteristic of the
time-inhomogeneous mean-field flow
\[
  \dot\theta(t)=\xi(t)\,b(\theta(t),\rho_t)
\]
started from $\supp(\rho_0)$ remains in the compact set
\[
  B_T^{\rm sgd}
  :=
  \{(a,c,\xi): |a|\le A_T^{\rm sgd},\ |c|\le C_T^{\rm sgd},\ \norm{\xi}\le R_T^{\rm sgd}\}.
\]
In particular, if
\[
  M_T^{\rm sgd}:=K_\varphi A_T^{\rm sgd},
\]
then along the deterministic path one has
\[
  |\eta_{1,\rho_t}(x)|\le M_T^{\rm sgd},
  \qquad
  \eta_{2,\rho_t}(x)\le -\lambda_{\min},
  \qquad
  \forall x\in\mc{X},\ \forall t\in[0,T].
\]
\end{proposition}

\begin{proof}
The proof is identical to Proposition~\ref{prop:softplus-confinement-k1}, with
the additional prefactor $\xi(t)$ absorbed by $\xi_{\max}$ in the comparison
argument.
\end{proof}

\begin{proposition}[Discrete-time finite-horizon confinement for SGD]
\label{prop:sgd-discrete-confinement-k1}
Assume Assumptions~\ref{ass:barrier-k1} and \ref{ass:softplus-sgd-k1}. Define
deterministic sequences $A_k^\varepsilon$, $C_k^\varepsilon$, and
$R_k^\varepsilon$ by
\begin{align}
  A_{k+1}^\varepsilon
  &=
  A_k^\varepsilon
  +
  s_k K_\varphi
  \left(
    \frac{K_\varphi A_k^\varepsilon}{2\lambda_{\min}}+Y_\ast
  \right),
  \qquad
  A_0^\varepsilon=A_0,
  \label{eq:discrete-A-envelope-k1}
  \\
  C_{k+1}^\varepsilon
  &=
  C_k^\varepsilon
  +
  s_k K_\varphi
  \left(
    \frac{K_\varphi^2 (A_k^\varepsilon)^2}{4\lambda_{\min}^2}
    +
    \frac{1}{2\lambda_{\min}}
    +
    Y_\ast^2
  \right),
  \qquad
  C_0^\varepsilon=C_0,
  \label{eq:discrete-C-envelope-k1}
  \\
  R_{k+1}^\varepsilon
  &=
  R_k^\varepsilon
  +
  s_k K_{\nabla}
  \Biggl[
    A_k^\varepsilon
    \left(
      \frac{K_\varphi A_k^\varepsilon}{2\lambda_{\min}}+Y_\ast
    \right)
    \nonumber\\
  &\qquad\qquad\qquad
    +
    C_k^\varepsilon
    \left(
      \frac{K_\varphi^2 (A_k^\varepsilon)^2}{4\lambda_{\min}^2}
      +
      \frac{1}{2\lambda_{\min}}
      +
      Y_\ast^2
    \right)
  \Biggr],
  \qquad
  R_0^\varepsilon=R_0.
  \label{eq:discrete-R-envelope-k1}
\end{align}
Then every SGD iterate of \eqref{eq:sgd-update-k1} satisfies
\[
  |a_i^k|\le A_k^\varepsilon,
  \qquad
  |c_i^k|\le C_k^\varepsilon,
  \qquad
  \norm{\xi_i^k}\le R_k^\varepsilon
  \qquad
  \forall i,\ \forall k\le T/\varepsilon.
\]
Moreover, the discrete envelopes are uniformly bounded by the continuous
majorants from Proposition~\ref{prop:sgd-confinement-k1}; in particular every
SGD iterate remains in $B_T^{\rm sgd}$.
\end{proposition}

\begin{proof}
We argue by induction on $k$. Assume
\[
  |a_i^k|\le A_k^\varepsilon,
  \qquad
  |c_i^k|\le C_k^\varepsilon,
  \qquad
  \norm{\xi_i^k}\le R_k^\varepsilon
\]
for all $i$. Then
\[
  |\eta_{1,\hat\rho_k^N}(x)|
  \le
  \frac1N\sum_{i=1}^N |a_i^k|\,|\varphi(\xi_i^k,x)|
  \le
  K_\varphi A_k^\varepsilon,
\]
and Proposition~\ref{prop:softplus-domain-k1} gives
$\eta_{2,\hat\rho_k^N}(x)\le -\lambda_{\min}$. Hence
\eqref{eq:q1-k1}--\eqref{eq:q2-k1} imply the same bounds as in the proof of
Proposition~\ref{prop:softplus-confinement-k1}, now with $A_k^\varepsilon$ in
place of $A_t$. Inserting these estimates into the single-sample update
\eqref{eq:sgd-update-k1} gives
\[
  |a_i^{k+1}|
  \le
  A_{k+1}^\varepsilon,
  \qquad
  |c_i^{k+1}|
  \le
  C_{k+1}^\varepsilon,
  \qquad
  \norm{\xi_i^{k+1}}
  \le
  R_{k+1}^\varepsilon.
\]
This closes the induction.

The comparison between the discrete recursion and the continuous majorant ODEs
is standard Euler-vs-ODE monotonicity: since $s_k\le \varepsilon\xi_{\max}$ and
the right-hand sides are nondecreasing in the envelopes, one has
$A_k^\varepsilon\le A_{k\varepsilon}^{\rm sgd}$, $C_k^\varepsilon\le
C_{k\varepsilon}^{\rm sgd}$, and $R_k^\varepsilon\le R_{k\varepsilon}^{\rm sgd}$
for all $k\le T/\varepsilon$.
\end{proof}

\begin{proposition}[Uniform stochastic-drift regularity]
\label{prop:sample-drift-regularity-k1}
Assume Assumptions~\ref{ass:barrier-k1}, \ref{ass:feature-second-k1}, and
\ref{ass:softplus-sgd-k1}. There exists
$\widehat L_T<\infty$ such that for every sample $z=(x,y)$ in the support of
$P$, all $\theta,\bar\theta\in B_T^{\rm sgd}$, and all probability measures
$\rho,\nu$ supported on $B_T^{\rm sgd}$,
\begin{align}
  \norm{\hat b(\theta,\rho;z)}
  &\le
  \widehat L_T,
  \label{eq:sample-drift-bound-k1}
  \\
  \norm{\hat b(\theta,\rho;z)-\hat b(\bar\theta,\nu;z)}
  &\le
  \widehat L_T\bigl(\norm{\theta-\bar\theta}+W_1(\rho,\nu)\bigr).
  \label{eq:sample-drift-lipschitz-k1}
\end{align}
\end{proposition}

\begin{proof}
On $B_T^{\rm sgd}$, the strip bound from
Proposition~\ref{prop:sgd-confinement-k1} and
Proposition~\ref{prop:gaussian-smoothness-k1} imply that
$q_1(\eta_\rho(x);y)$ and $q_2(\eta_\rho(x);y)$ are uniformly bounded and
Lipschitz in $(\theta,\rho)$ for $|y|\le Y_\ast$. Since $\psi'_{\rm sp}$ is
bounded and Lipschitz, $\varphi$ is $K_{\nabla}$-Lipschitz in $\xi$ by the
uniform gradient bound, and $\nabla_\xi\varphi$ is $L_{\nabla}$-Lipschitz in
$\xi$ by assumption, the formulas
\eqref{eq:sample-drift-a-k1}--\eqref{eq:sample-drift-xi-k1} immediately yield
\eqref{eq:sample-drift-bound-k1} and \eqref{eq:sample-drift-lipschitz-k1}.
\end{proof}

\begin{theorem}[Training-side theorem: finite-horizon SGD mean-field approximation]
\label{thm:sgd-mean-field-k1}
Assume Assumptions~\ref{ass:barrier-k1}, \ref{ass:feature-second-k1}, and
\ref{ass:softplus-sgd-k1}. Consider the noiseless SGD iterates
\eqref{eq:sgd-update-k1} and let $\rho_t$ be the deterministic measure-valued
path solving
\begin{equation}
  \partial_t \rho_t
  +
  \nabla_\theta\cdot\bigl(\rho_t\,\xi(t)b(\theta,\rho_t)\bigr)
  =
  0,
  \qquad
  \rho_{t=0}=\rho_0.
  \label{eq:sgd-pde-k1}
\end{equation}
Then:
\begin{enumerate}[label=(\roman*)]
  \item the PDE \eqref{eq:sgd-pde-k1} is well posed on $[0,T]$;
  \item there exists an error term $\delta_{N,\varepsilon}(T)\to 0$ as
  $N\to\infty$ and $\varepsilon\downarrow 0$ such that
  \[
    \sup_{k\le T/\varepsilon}
    W_1\!\bigl(\hat\rho_k^N,\rho_{k\varepsilon}\bigr)
    \le
    \delta_{N,\varepsilon}(T)
  \]
  with high probability;
  \item in bounded-support regimes, one obtains the same qualitative structure
  as in the Mei/Ferbach four-dynamics argument: a Monte-Carlo term of order
  $N^{-1/2}$ together with discretization and stochastic-gradient terms of
  order $\varepsilon$ and $\sqrt{\varepsilon}$.
\end{enumerate}
\end{theorem}

\begin{proof}[Proof sketch]
Propositions~\ref{prop:unbiased-sample-drift-k1},
\ref{prop:sgd-confinement-k1}, \ref{prop:sgd-discrete-confinement-k1}, and
\ref{prop:sample-drift-regularity-k1} verify the model-specific hypotheses
needed for the bounded-support four-dynamics comparison of
Mei--Misiakiewicz--Montanari~\cite{MeiMisiakiewiczMontanari2019}. Ferbach et
al.~\cite{Ferbach2024} use the same reduction in Appendix~C.2 for noiseless
regularization-free SGD. Applying that finite-horizon machinery to the present
raw stochastic drift \eqref{eq:sgd-update-k1} yields the stated approximation
theorem.
\end{proof}

\begin{corollary}[Training-side corollary: shared-law common path for SGD]
\label{cor:sgd-common-path-k1}
Assume Assumptions~\ref{ass:barrier-k1}, \ref{ass:feature-second-k1}, and
\ref{ass:softplus-sgd-k1}. Let $\hat\rho_{A,k}^N,\hat\rho_{B,k}^N$ be two
independent width-$N$ SGD runs of \eqref{eq:sgd-update-k1}, both initialized
i.i.d.\ from the same law $\rho_0$ and trained over the same horizon $[0,T]$.
Then
\[
  \sup_{k\le T/\varepsilon}
  W_1\!\bigl(\hat\rho_{A,k}^N,\hat\rho_{B,k}^N\bigr)
  \le
  2\,\delta_{N,\varepsilon}(T)
\]
with high probability. In particular,
\[
  \sup_{k\le T/\varepsilon}
  W_1\!\bigl(\hat\rho_{A,k}^N,\hat\rho_{B,k}^N\bigr)\to 0
\]
in probability as $N\to\infty$ and $\varepsilon\downarrow 0$.
\end{corollary}

\begin{proof}
Apply Theorem~\ref{thm:sgd-mean-field-k1} to each run and use the triangle
inequality.
\end{proof}

This shared-law common-path statement is the appendix-level form of the
training-side theorem used in the main text. The point is not a new abstract
SGD framework, but the fact that the practical natural-softplus Gaussian head
satisfies the bounded-support drift, confinement, and regularity hypotheses
required by the existing finite-horizon mean-field machinery.

\section{Interpolation-Side Barrier Proofs}

This appendix proves Theorem~\ref{thm:main-barrier-k1}. Exact convexity holds
in natural-parameter space, so the only remaining gap between parameter
interpolation and convex interpolation comes from vector-head transfer and
softplus curvature.
Throughout this appendix, we work under Assumption~\ref{ass:barrier-k1}.

\begin{lemma}[Automatic strip control along the interpolation path]
\label{lem:auto-strip-k1}
For all $x\in\mc{X}$ and all $t\in[0,1]$,
\[
  \norm{z_A(x)}_\infty,\ \norm{z_B(x)}_\infty,\ \norm{z_t(x)}_\infty,\
  \norm{\hat z_t(x)}_\infty
  \le
  K_\varphi M_V^\infty.
\]
Consequently,
\[
  \eta_{\Theta_A}(x),\ \eta_{\Theta_B}(x),\ \eta_{\Theta_t}(x),\
  \hat\eta_t(x),\ \tilde\eta_t(x)
  \in
  D_{K_\varphi M_V^\infty,\lambda_{\min}}
  \qquad
  \forall x,\ \forall t.
\]
\end{lemma}

\begin{proof}
For the first coordinate,
\[
  |u_{\Theta_A}(x)|
  \le
  \frac1N\sum_{i=1}^N |a_i^A|\,|\varphi(\xi_i^A,x)|
  \le
  K_\varphi
  \Bigl(\frac1N\sum_{i=1}^N (a_i^A)^2\Bigr)^{1/2}
  \le
  K_\varphi M_V^\infty.
\]
The same estimate holds for $s_{\Theta_A}$, and likewise for network $B$.
For the interpolated coefficients, convexity of the square gives
\[
  \frac1N\sum_{i=1}^N \bigl((1-t)a_i^A+ta_i^B\bigr)^2
  \le
  (1-t)\frac1N\sum_{i=1}^N (a_i^A)^2
  +
  t\frac1N\sum_{i=1}^N (a_i^B)^2
  \le
  (M_V^\infty)^2,
\]
and similarly for the $c$ coordinates. This proves the raw bound for $z_t$.
The same bound for $\hat z_t$ is immediate.

For the natural outputs, the first coordinate is unchanged and the second
coordinate is always $\le -\lambda_{\min}$ by
Proposition~\ref{prop:softplus-domain-k1}.
\end{proof}

\begin{proposition}[Softplus Jensen gap]
\label{prop:softplus-jensen-k1}
For any $s_A,s_B\in\R$ and $t\in[0,1]$,
\begin{equation}
  0
  \le
  (1-t)\psi_{\rm sp}(s_A)+t\psi_{\rm sp}(s_B)
  -
  \psi_{\rm sp}((1-t)s_A+ts_B)
  \le
  \frac18\,t(1-t)(s_A-s_B)^2.
  \label{eq:softplus-jensen-k1}
\end{equation}
Consequently, for every $x\in\mc{X}$,
\begin{equation}
  \norm{\tilde\eta_t(x)-\hat\eta_t(x)}_\infty
  \le
  \frac18\,t(1-t)\bigl(s_{\Theta_A}(x)-s_{\Theta_B}(x)\bigr)^2.
  \label{eq:natural-jensen-gap-k1}
\end{equation}
\end{proposition}

\begin{proof}
The function $\psi_{\rm sp}$ is convex and satisfies
\[
  \psi''_{\rm sp}(s)=\frac{e^s}{(1+e^s)^2}\le \frac14.
\]
For any $C^2$ convex function with second derivative bounded by $L$, the
interpolation error is at most $\frac{L}{2}t(1-t)(s_A-s_B)^2$. Applying this
with $L=\frac14$ gives \eqref{eq:softplus-jensen-k1}. Since only the second
natural coordinate is nonlinear, \eqref{eq:natural-jensen-gap-k1} is exactly
the same estimate written for $\tilde\eta_t$ and $\hat\eta_t$.
\end{proof}

\begin{proposition}[Direct vector-head transfer]
\label{prop:vector-head-transfer-k1}
For the aligned pair $(\Theta_A,\Theta_B)$,
\begin{equation}
  \sup_{t\in[0,1]}\ \sup_{x\in\mc{X}}
  \norm{z_t(x)-\hat z_t(x)}_\infty
  \le
  \frac12 K_{\nabla}M_V^\infty\sqrt{B_N}.
  \label{eq:vector-transfer-k1}
\end{equation}
In particular,
\[
  \sup_{t\in[0,1]}
  \Bigl(
    \E\big[\norm{z_t(X)-\hat z_t(X)}_\infty^2\big]
  \Bigr)^{1/2}
  \le
  \frac12 K_{\nabla}M_V^\infty\sqrt{B_N}.
\]
\end{proposition}

\begin{proof}
Set
\[
  \xi_i(t):=(1-t)\xi_i^A+t\xi_i^B,
  \qquad
  \Delta\xi_i:=\xi_i^A-\xi_i^B.
\]
Let $v_i^A=(a_i^A,c_i^A)$ and $v_i^B=(a_i^B,c_i^B)$. Then
\begin{align*}
  z_t(x)-\hat z_t(x)
  &=
  \frac1N\sum_{i=1}^N
  \Bigl[
    (1-t)v_i^A\bigl(\varphi(\xi_i(t),x)-\varphi(\xi_i^A,x)\bigr)
    \\
  &\qquad\qquad\qquad
    +
    tv_i^B\bigl(\varphi(\xi_i(t),x)-\varphi(\xi_i^B,x)\bigr)
  \Bigr].
\end{align*}
Since
\[
  \norm{\xi_i(t)-\xi_i^A}=t\norm{\Delta\xi_i},
  \qquad
  \norm{\xi_i(t)-\xi_i^B}=(1-t)\norm{\Delta\xi_i},
\]
the uniform gradient bound gives
\[
  |\varphi(\xi_i(t),x)-\varphi(\xi_i^A,x)|
  \le
  tK_{\nabla}\norm{\Delta\xi_i},
  \qquad
  |\varphi(\xi_i(t),x)-\varphi(\xi_i^B,x)|
  \le
  (1-t)K_{\nabla}\norm{\Delta\xi_i}.
\]
Taking the $j$th coordinate of the vector head and applying Cauchy--Schwarz,
\[
  |(z_t(x)-\hat z_t(x))_j|
  \le
  2t(1-t)K_{\nabla}M_V^\infty\sqrt{B_N}
  \le
  \frac12 K_{\nabla}M_V^\infty\sqrt{B_N}.
\]
Taking the maximum over $j\in\{1,2\}$ proves the claim.
\end{proof}

\begin{proposition}[Natural-output interpolation gap]
\label{prop:natural-output-gap-k1}
For every $t\in[0,1]$,
\begin{equation}
  \E\big[\norm{\eta_{\Theta_t}(X)-\hat\eta_t(X)}_\infty\big]
  \le
  \frac12 K_{\nabla}M_V^\infty\sqrt{B_N}
  +
  \frac18\,t(1-t)\Delta_{s,N}.
  \label{eq:natural-gap-k1}
\end{equation}
\end{proposition}

\begin{proof}
By the triangle inequality,
\[
  \norm{\eta_{\Theta_t}(x)-\hat\eta_t(x)}_\infty
  \le
  \norm{\eta_{\Theta_t}(x)-\tilde\eta_t(x)}_\infty
  +
  \norm{\tilde\eta_t(x)-\hat\eta_t(x)}_\infty.
\]
The map $\Gamma(u,s)=(u,-\psi_{\rm sp}(s))$ is $1$-Lipschitz in the
$\ell_\infty$ norm because its first coordinate is the identity and
$0<\psi'_{\rm sp}(s)<1$. Hence
\[
  \norm{\eta_{\Theta_t}(x)-\tilde\eta_t(x)}_\infty
  \le
  \norm{z_t(x)-\hat z_t(x)}_\infty,
\]
and Proposition~\ref{prop:vector-head-transfer-k1} bounds the expectation of
this term by $\frac12 K_{\nabla}M_V^\infty\sqrt{B_N}$. The second term is
controlled by Proposition~\ref{prop:softplus-jensen-k1}.
\end{proof}

\begin{theorem}[Deterministic aligned barrier theorem]
\label{thm:deterministic-lmc-k1}
Assume Assumption~\ref{ass:barrier-k1}. For every fixed neuron permutation
$\varpi$,
\begin{equation}
  \cB(\Theta_A,\Theta_B;\varpi)
  \le
  C_\eta(K_\varphi M_V^\infty,\lambda_{\min},Y_\ast)
  \left[
    \frac12 K_{\nabla}M_V^\infty\sqrt{B_N}
    +
    \frac1{32}\Delta_{s,N}
  \right].
  \label{eq:deterministic-barrier-k1}
\end{equation}
\end{theorem}

\begin{proof}
Fix $t\in[0,1]$. By Lemma~\ref{lem:auto-strip-k1}, all relevant outputs lie in
$D_{K_\varphi M_V^\infty,\lambda_{\min}}$, so
Proposition~\ref{prop:convex-lipschitz-strip-k1} applies there. First use the
Lipschitz bound:
\[
  \E[\ell_\eta(\eta_{\Theta_t}(X);Y)]
  \le
  \E[\ell_\eta(\hat\eta_t(X);Y)]
  +
  C_\eta(K_\varphi M_V^\infty,\lambda_{\min},Y_\ast)
  \E[\norm{\eta_{\Theta_t}(X)-\hat\eta_t(X)}_\infty].
\]
Now apply Proposition~\ref{prop:natural-output-gap-k1}. The remaining term is
handled by convexity:
\[
  \ell_\eta(\hat\eta_t(X);Y)
  \le
  (1-t)\ell_\eta(\eta_{\Theta_A}(X);Y)
  +
  t\ell_\eta(\eta_{\Theta_B}(X);Y).
\]
Taking expectations and then $\sup_{t\in[0,1]}$ gives
\eqref{eq:deterministic-barrier-k1}, since $\sup_t t(1-t)=1/4$.
\end{proof}

\section{Wasserstein Alignment and LMC Proofs}

This appendix proves Theorem~\ref{thm:main-sgd-lmc-k1}. The alignment step
follows the optimal-transport route of Ferbach et al.~\cite{Ferbach2024}, while
the new model-specific ingredient is the natural-parameter barrier theorem from
Appendix~B.

\begin{lemma}[Permutation selected by empirical Wasserstein distance]
\label{lem:w1-permutation-k1}
Let
\[
  \rho_A^N:=\frac1N\sum_{i=1}^N\delta_{\theta_i^A},
  \qquad
  \rho_B^N:=\frac1N\sum_{i=1}^N\delta_{\theta_i^B},
  \qquad
  \theta_i^X:=(a_i^X,c_i^X,\xi_i^X).
\]
Assume that for some deterministic bounds $A,C,R<\infty$,
\[
  |a_i^A|,\ |a_i^B|\le A,
  \qquad
  |c_i^A|,\ |c_i^B|\le C,
  \qquad
  \norm{\xi_i^A},\ \norm{\xi_i^B}\le R
  \qquad
  \forall i=1,\dots,N.
\]
Let $\varepsilon_N:=W_1(\rho_A^N,\rho_B^N)$, where $W_1$ is computed with the
Euclidean cost on $\R^{2+d}$. Then there exists a permutation $\varpi_N$ such
that
\begin{equation}
  \frac1N\sum_{i=1}^N\norm{\theta_i^A-\theta_{\varpi_N(i)}^B}
  =
  \varepsilon_N.
  \label{eq:w1-permutation-k1}
\end{equation}
After relabeling network $B$ by this permutation, the alignment quantities
satisfy
\begin{align}
  B_N
  &\le
  2R\,\varepsilon_N,
  \label{eq:w1-to-BN-k1}
  \\
  \Delta_{s,N}
  &\le
  \bigl(K_\varphi+CK_{\nabla}\bigr)^2\varepsilon_N^2,
  \label{eq:w1-to-scale-k1}
\end{align}
\end{lemma}

\begin{proof}
Because both empirical measures place mass $1/N$ at each atom, the optimal
transport plan for $W_1$ is realized by a permutation, proving
\eqref{eq:w1-permutation-k1}.

For \eqref{eq:w1-to-BN-k1}, note that
\[
  \norm{\xi_i^A-\xi_i^B}\le 2R
  \qquad\text{and}\qquad
  \norm{\xi_i^A-\xi_i^B}\le \norm{\theta_i^A-\theta_i^B}.
\]
Hence
\[
  \norm{\xi_i^A-\xi_i^B}^2
  \le
  2R\norm{\theta_i^A-\theta_i^B}.
\]
Averaging proves \eqref{eq:w1-to-BN-k1}.

For \eqref{eq:w1-to-scale-k1}, write
\begin{align*}
  s_{\Theta_A}(x)-s_{\Theta_B}(x)
  &=
  \frac1N\sum_{i=1}^N
  (c_i^A-c_i^B)\varphi(\xi_i^A,x)
  \\
  &\qquad
  +
  \frac1N\sum_{i=1}^N
  c_i^B\bigl(\varphi(\xi_i^A,x)-\varphi(\xi_i^B,x)\bigr).
\end{align*}
Using $|\varphi|\le K_\varphi$, $|c_i^B|\le C$, and the gradient bound on
$\varphi$,
\[
  |s_{\Theta_A}(x)-s_{\Theta_B}(x)|
  \le
  \bigl(K_\varphi+CK_{\nabla}\bigr)\varepsilon_N.
\]
Squaring and taking expectation in $X$ yields \eqref{eq:w1-to-scale-k1}.
\end{proof}

\begin{theorem}[LMC from a common path]
\label{thm:lmc-from-common-path-k1}
Assume the trained endpoints satisfy the deterministic bounds
\[
  |a_i^A|,\ |a_i^B|\le A_T,
  \qquad
  |c_i^A|,\ |c_i^B|\le C_T,
  \qquad
  \norm{\xi_i^A},\ \norm{\xi_i^B}\le R_T
  \qquad
  \forall i=1,\dots,N,
\]
and let
\[
  \varepsilon_N:=W_1(\rho_A^N,\rho_B^N).
\]
Let $\varpi_N$ be the permutation from Lemma~\ref{lem:w1-permutation-k1}. Then
\begin{equation}
  \cB(\Theta_A,\Theta_B;\varpi_N)
  \le
  C_\eta(K_\varphi M_{V,T},\lambda_{\min},Y_\ast)
  \left[
    \frac12 K_{\nabla}M_{V,T}\sqrt{2R_T}\,\varepsilon_N^{1/2}
    +
    \frac{(K_\varphi+C_TK_{\nabla})^2}{32}\,\varepsilon_N^2
  \right],
  \label{eq:lmc-common-path-bound-k1}
\end{equation}
In particular, if $\varepsilon_N\to 0$, then
\[
  \cB(\Theta_A,\Theta_B;\varpi_N)\to 0
\]
in the same mode of convergence.
\end{theorem}

\begin{proof}
The displayed deterministic bounds imply $M_V^\infty\le M_{V,T}$. Apply
Theorem~\ref{thm:deterministic-lmc-k1} after relabeling network $B$ by
$\varpi_N$, and then use Lemma~\ref{lem:w1-permutation-k1} with
$(A,C,R)=(A_T,C_T,R_T)$ together with \eqref{eq:w1-to-BN-k1} and
\eqref{eq:w1-to-scale-k1} to
replace $B_N$ and $\Delta_{s,N}$ by Wasserstein-based quantities.
\end{proof}

\section{Rate Refinements}

This appendix collects secondary rate statements. They are not needed for the
main finite-horizon common-path to LMC chain in the body of the paper.
Throughout the appendix, assume the aligned endpoints satisfy the deterministic
bounds of Theorem~\ref{thm:lmc-from-common-path-k1}, and write
\[
  M_{V,T}:=\max\{A_T,C_T\}.
\]

\subsection{Baseline rate from hidden-law concentration}

\begin{proposition}[From hidden-law concentration to $B_N=O(d/N)$]
\label{prop:BN-from-w2-k1}
Let
\[
  \hat\nu_A^N:=\frac1N\sum_{i=1}^N \delta_{\xi_i^A},
  \qquad
  \hat\nu_B^N:=\frac1N\sum_{i=1}^N \delta_{\xi_i^B},
\]
and let $\mu_{\xi,\ast}$ be a reference law on $\R^d$. Suppose that for some
deterministic sequence $\eta_N\downarrow 0$,
\[
  W_2(\hat\nu_A^N,\mu_{\xi,\ast})\le \eta_N,
  \qquad
  W_2(\hat\nu_B^N,\mu_{\xi,\ast})\le \eta_N.
\]
If the neuron permutation $\varpi_N$ is chosen by hidden-weight optimal transport,
then
\begin{equation}
  B_N = W_2^2(\hat\nu_A^N,\hat\nu_B^N)\le 4\eta_N^2.
  \label{eq:BN-via-w2-k1}
\end{equation}
In particular, if
\[
  \eta_N=O(\kappa_\xi\sqrt{d/N}),
\]
then
\[
  B_N=O(\kappa_\xi^2 d/N).
\]
\end{proposition}

\begin{proof}
Because the hidden empirical laws have equal weights $1/N$, the $W_2$-optimal
coupling between them is realized by a permutation, and the corresponding
squared cost is exactly $B_N$. The triangle inequality gives
\[
  W_2(\hat\nu_A^N,\hat\nu_B^N)
  \le
  W_2(\hat\nu_A^N,\mu_{\xi,\ast}) + W_2(\hat\nu_B^N,\mu_{\xi,\ast})
  \le
  2\eta_N.
\]
Squaring proves \eqref{eq:BN-via-w2-k1}.
\end{proof}

\begin{corollary}[Baseline rate]
\label{cor:baseline-rate-k1}
Suppose that along a width sequence,
\[
  B_N=O(d/N),
  \qquad
  \Delta_{s,N}=O(1/N).
\]
Then Theorem~\ref{thm:deterministic-lmc-k1} yields
\[
  \cB(\Theta_A,\Theta_B;\varpi)
  =
  O\!\left(\sqrt{\frac{d}{N}}+\frac1N\right).
\]
In particular, for fixed input dimension $d$,
\[
  \cB(\Theta_A,\Theta_B;\varpi)=O(N^{-1/2}).
\]
\end{corollary}

\begin{proof}
Substitute the displayed rates into \eqref{eq:deterministic-barrier-k1}.
\end{proof}

\subsection{Conditional head regularity and refined transfer}

A faster route is available if the trained head varies regularly with the
hidden feature, improving the transfer term from $O(\sqrt{B_N})$ to $O(B_N)$.

\begin{definition}[Per-coordinate head mismatch]
\label{def:head-mismatch-k1}
Define
\[
  M_{\Delta V}^\infty
  :=
  \max_{j\in\{1,2\}}
  \Bigl(
    \frac1N\sum_{i=1}^N \bigl((v_i^A-v_i^B)_j\bigr)^2
  \Bigr)^{1/2},
  \qquad
  v_i^X:=(a_i^X,c_i^X)\in\R^2.
\]
\end{definition}

\begin{assumption}[Conditional head regularity]
\label{ass:head-regularity-k1}
There exist a map $g:\R^d\to\R^2$ and residuals $r_i^A,r_i^B\in\R^2$ such that
\[
  v_i^A=g(\xi_i^A)+r_i^A,
  \qquad
  v_i^B=g(\xi_i^B)+r_i^B,
\]
and
\[
  \norm{g(\xi)-g(\bar\xi)}_\infty
  \le
  L_g\norm{\xi-\bar\xi}
  \qquad
  \forall \xi,\bar\xi\in\R^d.
\]
Define the residual mismatch
\[
  \mc{R}_N
  :=
  \max_{j\in\{1,2\}}
  \Bigl(
    \frac1N\sum_{i=1}^N \bigl((r_i^A-r_i^B)_j\bigr)^2
  \Bigr)^{1/2}.
\]
\end{assumption}

\begin{lemma}[Head mismatch from conditional regularity]
\label{lem:head-mismatch-regular-k1}
Under Assumption~\ref{ass:head-regularity-k1},
\begin{equation}
  M_{\Delta V}^\infty
  \le
  L_g\sqrt{B_N}+\mc{R}_N.
  \label{eq:head-mismatch-regular-k1}
\end{equation}
\end{lemma}

\begin{proof}
For each coordinate $j\in\{1,2\}$,
\[
  |(v_i^A-v_i^B)_j|
  \le
  |(g(\xi_i^A)-g(\xi_i^B))_j|
  +
  |(r_i^A-r_i^B)_j|
  \le
  L_g\norm{\xi_i^A-\xi_i^B}+|(r_i^A-r_i^B)_j|.
\]
Take the empirical $\ell_2$ norm in $i$ and the maximum over $j$.
\end{proof}

\begin{proposition}[Refined vector-head transfer]
\label{prop:refined-transfer-k1}
Under Assumptions~\ref{ass:barrier-k1} and \ref{ass:feature-second-k1}, one has
\begin{equation}
  \sup_{t\in[0,1]}\ \sup_{x\in\mc{X}}
  \norm{z_t(x)-\hat z_t(x)}_\infty
  \le
  \frac{K_{\nabla}}{4}M_{\Delta V}^\infty\sqrt{B_N}
  +
  \frac{L_{\nabla}}{8}M_{V,T}B_N,
  \label{eq:refined-transfer-k1}
\end{equation}
where $M_{V,T}:=\max\{A_T,C_T\}$.
\end{proposition}

\begin{proof}
Write
\[
  \phi_i^A(x):=\varphi(\xi_i^A,x),
  \qquad
  \phi_i^B(x):=\varphi(\xi_i^B,x),
  \qquad
  \phi_i^t(x):=\varphi((1-t)\xi_i^A+t\xi_i^B,x).
\]
Then
\begin{align*}
  z_t(x)-\hat z_t(x)
  &=
  \frac1N\sum_{i=1}^N
  \Bigl(
    v_i^t\phi_i^t-(1-t)v_i^A\phi_i^A-tv_i^B\phi_i^B
  \Bigr)
  \\
  &=
  \frac1N\sum_{i=1}^N
  v_i^t\bigl(\phi_i^t-(1-t)\phi_i^A-t\phi_i^B\bigr)
  \\
  &\qquad
  +
  t(1-t)\frac1N\sum_{i=1}^N
  (v_i^B-v_i^A)(\phi_i^A-\phi_i^B),
\end{align*}
where $v_i^t:=(1-t)v_i^A+tv_i^B$.

For the second sum, use
\[
  |\phi_i^A(x)-\phi_i^B(x)|
  \le
  K_{\nabla}\norm{\xi_i^A-\xi_i^B}
\]
and Cauchy--Schwarz to obtain
\[
  t(1-t)\frac1N\sum_{i=1}^N
  |(v_i^B-v_i^A)_j|\,|\phi_i^A(x)-\phi_i^B(x)|
  \le
  \frac{K_{\nabla}}{4}M_{\Delta V}^\infty\sqrt{B_N}.
\]

For the first sum, the Lipschitz continuity of $\nabla_\xi\varphi$ gives the
second-order interpolation remainder
\[
  \bigl|
    \phi_i^t(x)-(1-t)\phi_i^A(x)-t\phi_i^B(x)
  \bigr|
  \le
  \frac{L_{\nabla}}{2}\,t(1-t)\norm{\xi_i^A-\xi_i^B}^2.
\]
Since $|(v_i^t)_j|\le M_{V,T}$ under the finite-horizon confinement bound,
\[
  \frac1N\sum_{i=1}^N
  |(v_i^t)_j|
  \bigl|
    \phi_i^t(x)-(1-t)\phi_i^A(x)-t\phi_i^B(x)
  \bigr|
  \le
  \frac{L_{\nabla}}{8}M_{V,T}B_N.
\]
Adding the two estimates and taking the maximum over $j$ proves the claim.
\end{proof}

\begin{proposition}[Endpoint raw-scale rate]
\label{prop:scale-rate-k1}
Under the standing appendix bounds,
\begin{equation}
  \Delta_{s,N}^{1/2}
  \le
  K_\varphi M_{\Delta V}^\infty + C_TK_{\nabla}\sqrt{B_N}.
  \label{eq:scale-rate-k1}
\end{equation}
Consequently, if
\[
  M_{\Delta V}^\infty=O(\sqrt{B_N}),
\]
then
\[
  \Delta_{s,N}=O(B_N).
\]
\end{proposition}

\begin{proof}
Exactly as in the proof of Lemma~\ref{lem:w1-permutation-k1},
\begin{align*}
  s_{\Theta_A}(x)-s_{\Theta_B}(x)
  &=
  \frac1N\sum_{i=1}^N
  (c_i^A-c_i^B)\varphi(\xi_i^A,x)
  \\
  &\qquad
  +
  \frac1N\sum_{i=1}^N
  c_i^B\bigl(\varphi(\xi_i^A,x)-\varphi(\xi_i^B,x)\bigr).
\end{align*}
The first term is bounded by $K_\varphi M_{\Delta V}^\infty$ and the second by
$C_TK_{\nabla}\sqrt{B_N}$. Take the supremum in $x$ and then the
$L^2(P_X)$ norm.
\end{proof}

\begin{corollary}[Fast rate under conditional head regularity]
\label{cor:fast-rate-k1}
Assume Assumptions~\ref{ass:barrier-k1}, \ref{ass:feature-second-k1}, and
\ref{ass:head-regularity-k1}. If
\[
  \mc{R}_N=O(\sqrt{B_N}).
\]
Then
\[
  \cB(\Theta_A,\Theta_B;\varpi)=O(B_N).
\]
In particular, if
\[
  B_N=O(d/N),
\]
then
\[
  \cB(\Theta_A,\Theta_B;\varpi)=O(d/N),
\]
which is $O(1/N)$ for fixed input dimension $d$.
\end{corollary}

\begin{proof}
Lemma~\ref{lem:head-mismatch-regular-k1} implies
$M_{\Delta V}^\infty=O(\sqrt{B_N})$. Proposition~\ref{prop:refined-transfer-k1}
then yields
\[
  \sup_t\sup_x \norm{z_t(x)-\hat z_t(x)}_\infty = O(B_N).
\]
Because $\Gamma$ is $1$-Lipschitz, the same bound holds for
$\eta_{\Theta_t}-\tilde\eta_t$. Proposition~\ref{prop:scale-rate-k1} gives
$\Delta_{s,N}=O(B_N)$, hence Proposition~\ref{prop:softplus-jensen-k1} gives
\[
  \sup_t \E\big[\norm{\tilde\eta_t(X)-\hat\eta_t(X)}_\infty\big]=O(B_N).
\]
Repeating the proof of Theorem~\ref{thm:deterministic-lmc-k1} with these
refined estimates gives $\cB=O(B_N)$.
\end{proof}

\section{Sharpened Deterministic Barrier Refinements}

This appendix sharpens the deterministic interpolation-side barrier theorem of
Appendix~B. The baseline route passes through a global natural strip and a
single natural-space Lipschitz constant. The refinements below replace that
last step by a raw-space reduction, exact sigmoid-projected transfer
functionals, and exact conditional-risk moduli. These statements are
deterministic appendix-level refinements: they are not needed for the main
finite-horizon LMC theorem, but they materially tighten theorem-consistent
instantiated certificates on the matched softplus run.

\subsection{Raw-Space Barrier Reduction}

Write the Gaussian natural loss as
\[
  \ell_\eta(\eta;y)=A_{\rm G}(\eta)-\eta_1 y-\eta_2 y^2,
  \qquad
  \eta=(\eta_1,\eta_2),\ \eta_2<0.
\]
On the natural strip
\[
  D_{M,\lambda}
  :=
  \{(\eta_1,\eta_2)\in\R^2:\ |\eta_1|\le M,\ \eta_2\le-\lambda\},
\]
the second-coordinate loss constant is
\begin{equation}
  C_2(M,\lambda,Y_\ast)
  :=
  \sup_{D_{M,\lambda},\,|y|\le Y_\ast}
  \abs{\partial_{\eta_2}\ell_\eta(\eta;y)}
  \le
  \frac{M^2}{4\lambda^2}+\frac{1}{2\lambda}+Y_\ast^2.
  \label{eq:C2}
\end{equation}

Now work in raw coordinates
\[
  z=(u,s),
  \qquad
  \tilde\ell(u,s;y):=\ell_\eta(\Gamma(u,s);y).
\]
Define the raw strip
\[
  \mathcal S_{M,\lambda}
  :=
  \{(u,s)\in\R^2:\ |u|\le M,\ \psi_{\rm sp}(s)\ge \lambda\}.
\]

\begin{proposition}[Raw-strip derivative constants]
\label{prop:raw-strip-constants}
Assume $\lambda\ge \lambda_{\min}$. On $\mathcal S_{M,\lambda}$ and for
$|y|\le Y_\ast$, define
\begin{align}
  C_u^{\rm raw}(M,\lambda,Y_\ast)
  &:=
  \sup_{\mathcal S_{M,\lambda},\,|y|\le Y_\ast}
  \abs{\partial_u\tilde\ell(u,s;y)}
  \le
  \frac{M}{2\lambda}+Y_\ast,
  \label{eq:Cu-raw}
  \\
  C_s^{\rm raw}(M,\lambda,Y_\ast)
  &:=
  \sup_{\mathcal S_{M,\lambda},\,|y|\le Y_\ast}
  \abs{\partial_s\tilde\ell(u,s;y)}
  \nonumber\\
  &\le
  \sup_{r\ge \lambda}
  \Bigl(1-e^{-(r-\lambda_{\min})}\Bigr)
  \left(
    \frac{M^2}{4r^2}+\frac1{2r}+Y_\ast^2
  \right).
  \label{eq:Cs-raw}
\end{align}
\end{proposition}

\begin{proof}
Write $r=\psi_{\rm sp}(s)$. Then
\[
  \partial_u\tilde\ell(u,s;y)=\frac{u}{2r}-y,
  \qquad
  \partial_s\tilde\ell(u,s;y)
  =
  -\psi_{\rm sp}'(s)\,\partial_{\eta_2}\ell_\eta(\Gamma(u,s);y).
\]
Since
\[
  \psi_{\rm sp}'(s)
  =
  \frac{e^s}{1+e^s}
  =
  1-e^{-(r-\lambda_{\min})},
\]
and
\[
  \abs{\partial_{\eta_2}\ell_\eta(\Gamma(u,s);y)}
  \le
  \frac{M^2}{4r^2}+\frac1{2r}+Y_\ast^2,
\]
the claim follows.
\end{proof}

\begin{proposition}[Raw-space barrier reduction]
\label{prop:rawspace-barrier}
Assume $\lambda\ge \lambda_{\min}$ and that, for every $t\in[0,1]$ and every
$x\in\mc{X}$, both
$z_t(x)=(u_t(x),s_t(x))$ and $\hat z_t(x)=(\hat u_t(x),\hat s_t(x))$ lie in
$\mathcal S_{M,\lambda}$. Define
\[
  U:=\sup_{t\in[0,1]}\E\big[\abs{u_t(X)-\hat u_t(X)}\big],
  \qquad
  S:=\sup_{t\in[0,1]}\E\big[\abs{s_t(X)-\hat s_t(X)}\big].
\]
Then
\begin{equation}
  \cB(\Theta_A,\Theta_B;\varpi)
  \le
  C_u^{\rm raw}(M,\lambda,Y_\ast)\,U
  +
  C_s^{\rm raw}(M,\lambda,Y_\ast)\,S
  +
  C_2(M,\lambda,Y_\ast)\frac1{32}\Delta_{s,N}.
  \label{eq:rawspace-barrier}
\end{equation}
\end{proposition}

\begin{proof}
Convexity in natural space gives
\[
  \cB(\Theta_A,\Theta_B;\varpi)
  \le
  \sup_t
  \E\big[
    \ell_\eta(\eta_{\Theta_t}(X);Y)-\ell_\eta(\hat\eta_t(X);Y)
  \big].
\]
Insert the intermediate output
\[
  \tilde\eta_t(X)=\Gamma(\hat z_t(X)),
  \qquad
  \ell_\eta(\tilde\eta_t(X);Y)=\tilde\ell(\hat z_t(X);Y).
\]
Then
\begin{align*}
  \ell_\eta(\eta_{\Theta_t}(X);Y)-\ell_\eta(\hat\eta_t(X);Y)
  &=
  \bigl[
    \tilde\ell(z_t(X);Y)-\tilde\ell(\hat z_t(X);Y)
  \bigr]
  \\
  &\quad+
  \bigl[
    \ell_\eta(\tilde\eta_t(X);Y)-\ell_\eta(\hat\eta_t(X);Y)
  \bigr].
\end{align*}
On $\mathcal S_{M,\lambda}$, the first bracket is bounded by the mean value
theorem:
\[
  C_u^{\rm raw}(M,\lambda,Y_\ast)\abs{u_t(X)-\hat u_t(X)}
  +
  C_s^{\rm raw}(M,\lambda,Y_\ast)\abs{s_t(X)-\hat s_t(X)}.
\]
For the second bracket, $\tilde\eta_t(X)$ and $\hat\eta_t(X)$ have the same
first coordinate, so the mean value theorem in the $\eta_2$ direction gives
\begin{align*}
  \abs{
    \ell_\eta(\tilde\eta_t(X);Y)-\ell_\eta(\hat\eta_t(X);Y)
  }
  &\le
  C_2(M,\lambda,Y_\ast)\,
  \Bigl|
    \psi_{\rm sp}(\hat s_t(X))
    -(1-t)\psi_{\rm sp}(s_{\Theta_A}(X))
  \\
  &\qquad\qquad
    -t\psi_{\rm sp}(s_{\Theta_B}(X))
  \Bigr|.
\end{align*}
The expectation of the displayed softplus Jensen defect is bounded by
$\frac1{32}\Delta_{s,N}$ by Proposition~\ref{prop:softplus-jensen-k1}. Taking
expectations and $\sup_t$ yields \eqref{eq:rawspace-barrier}.
\end{proof}

\subsection{Exact Sigmoid-Projected Transfer}

Specialize to the experimental hidden feature
\[
  \varphi(\xi,x)=\sigma(\xi_1x+\xi_0),
  \qquad
  \sigma(z)=(1+e^{-z})^{-1}.
\]
For matched endpoint parameters $\xi_i^A,\xi_i^B\in\R^2$, write
\[
  z_i^A(x):=\langle \xi_i^A,\tilde x\rangle,
  \qquad
  z_i^B(x):=\langle \xi_i^B,\tilde x\rangle,
  \qquad
  z_{i,t}(x):=(1-t)z_i^A(x)+t z_i^B(x),
  \qquad
  \tilde x:=(x,1).
\]

\begin{definition}[Exact sigmoid-projected transfer functionals]
\label{def:exact-sigmoid-projected}
Define the averaged transfer functionals
\begin{align}
  d_i
  &:=
  \E\big[\abs{\sigma(z_i^A(X))-\sigma(z_i^B(X))}\big],
  \label{eq:exact-di}
  \\
  r_i^{(a)}
  &:=
  \sup_{t\in[0,1]}
  \abs{(1-t)a_i^A+t a_i^B}\,
  \E\big[
    \abs{\sigma(z_{i,t}(X))-(1-t)\sigma(z_i^A(X))-t\sigma(z_i^B(X))}
  \big],
  \label{eq:exact-ra}
  \\
  r_i^{(c)}
  &:=
  \sup_{t\in[0,1]}
  \abs{(1-t)c_i^A+t c_i^B}\,
  \E\big[
    \abs{\sigma(z_{i,t}(X))-(1-t)\sigma(z_i^A(X))-t\sigma(z_i^B(X))}
  \big].
  \label{eq:exact-rc}
\end{align}
Likewise define the pointwise versions
\begin{align}
  d_i^\infty
  &:=
  \sup_{x\in\mc{X}}
  \abs{\sigma(z_i^A(x))-\sigma(z_i^B(x))},
  \label{eq:exact-di-inf}
  \\
  r_i^{(a,\infty)}
  &:=
  \sup_{t\in[0,1]}\sup_{x\in\mc{X}}
  \abs{(1-t)a_i^A+t a_i^B}\,
  \abs{\sigma(z_{i,t}(x))-(1-t)\sigma(z_i^A(x))-t\sigma(z_i^B(x))},
  \label{eq:exact-ra-inf}
  \\
  r_i^{(c,\infty)}
  &:=
  \sup_{t\in[0,1]}\sup_{x\in\mc{X}}
  \abs{(1-t)c_i^A+t c_i^B}\,
  \abs{\sigma(z_{i,t}(x))-(1-t)\sigma(z_i^A(x))-t\sigma(z_i^B(x))}.
  \label{eq:exact-rc-inf}
\end{align}
\end{definition}

\begin{proposition}[Exact sigmoid-projected refined transfer]
\label{prop:exact-sigmoid-projected-transfer}
Under Definition~\ref{def:exact-sigmoid-projected},
\begin{align}
  \sup_{t\in[0,1]}\E\big[\abs{u_t(X)-\hat u_t(X)}\big]
  &\le
  \frac{1}{4N}\sum_{i=1}^N \abs{a_i^A-a_i^B}\,d_i
  +
  \frac{1}{N}\sum_{i=1}^N r_i^{(a)},
  \label{eq:exact-transfer-u}
  \\
  \sup_{t\in[0,1]}\E\big[\abs{s_t(X)-\hat s_t(X)}\big]
  &\le
  \frac{1}{4N}\sum_{i=1}^N \abs{c_i^A-c_i^B}\,d_i
  +
  \frac{1}{N}\sum_{i=1}^N r_i^{(c)},
  \label{eq:exact-transfer-s}
\end{align}
and
\begin{align}
  \sup_{t\in[0,1]}\sup_{x\in\mc{X}}\abs{u_t(x)-\hat u_t(x)}
  &\le
  \frac{1}{4N}\sum_{i=1}^N \abs{a_i^A-a_i^B}\,d_i^\infty
  +
  \frac{1}{N}\sum_{i=1}^N r_i^{(a,\infty)},
  \label{eq:exact-transfer-u-inf}
  \\
  \sup_{t\in[0,1]}\sup_{x\in\mc{X}}\abs{s_t(x)-\hat s_t(x)}
  &\le
  \frac{1}{4N}\sum_{i=1}^N \abs{c_i^A-c_i^B}\,d_i^\infty
  +
  \frac{1}{N}\sum_{i=1}^N r_i^{(c,\infty)}.
  \label{eq:exact-transfer-s-inf}
\end{align}
\end{proposition}

\begin{proof}
Repeat the refined transfer decomposition neuron by neuron, but keep the exact
feature difference and the exact interpolation defect instead of replacing them
by global bounds on $\sigma'$ and $\sigma''$. The mixed head-difference term
still carries the factor $t(1-t)\le 1/4$, while the interpolation remainder is
absorbed directly into the functionals
\eqref{eq:exact-ra}--\eqref{eq:exact-rc-inf}.
\end{proof}

\subsection{Endpoint Variance Strip}

\begin{proposition}[Pointwise conditional risk implies an endpoint variance strip]
\label{prop:conditional-risk-strip}
For a network $\Theta$, define the conditional risk
\[
  r_\Theta(x):=\E\big[\ell_\eta(\eta_\Theta(x);Y)\mid X=x\big].
\]
If
\[
  r_\Theta(x)\le L_{\rm end}
  \qquad
  \forall x\in\mc{X},
\]
then
\[
  -\eta_{2,\Theta}(x)\ge \pi e^{-2L_{\rm end}}
  \qquad
  \forall x\in\mc{X}.
\]
\end{proposition}

\begin{proof}
In mean/variance coordinates,
\[
  \ell_\eta(\eta_\Theta(x);Y)
  =
  \frac12\log(2\pi v_\Theta(x))
  +
  \frac{(Y-\mu_\Theta(x))^2}{2v_\Theta(x)},
  \qquad
  v_\Theta(x)=\frac{1}{2(-\eta_{2,\Theta}(x))}.
\]
Drop the nonnegative quadratic term to get
\[
  \frac12\log(2\pi v_\Theta(x))\le L_{\rm end}.
\]
Exponentiating yields
\[
  v_\Theta(x)\le \frac{e^{2L_{\rm end}}}{2\pi},
\]
hence
\[
  -\eta_{2,\Theta}(x)=\frac{1}{2v_\Theta(x)}\ge \pi e^{-2L_{\rm end}}.
\]
\end{proof}

\begin{remark}
Conditional risk sharpens the variance strip but does not control $\eta_1$.
This is exactly why the raw-space route is preferable: it controls the first
coordinate by the architectural raw bound on $u$, and uses conditional risk
only on the variance side.
\end{remark}

\begin{proposition}[Raw-scale propagation of an endpoint variance strip]
\label{prop:raw-scale-strip-propagation}
Assume the endpoints satisfy
\[
  -\eta_{2,\Theta_A}(x),\ -\eta_{2,\Theta_B}(x)\ge \lambda_{\rm end}
  \qquad
  \forall x\in\mc{X},
\]
for some $\lambda_{\rm end}>\lambda_{\min}$. Let
\[
  \beta_{\rm end}:=\log\bigl(e^{\lambda_{\rm end}-\lambda_{\min}}-1\bigr),
  \qquad
  \delta_s:=\sup_{t\in[0,1]}\sup_{x\in\mc{X}}\abs{s_t(x)-\hat s_t(x)}.
\]
Then
\[
  -\eta_{2,\Theta_t}(x)\ge \psi_{\rm sp}(\beta_{\rm end}-\delta_s)
  \qquad
  \forall t\in[0,1],\ \forall x\in\mc{X}.
\]
\end{proposition}

\begin{proof}
The endpoint assumption is equivalent to
\[
  s_{\Theta_A}(x),\ s_{\Theta_B}(x)\ge \beta_{\rm end},
\]
hence their raw interpolation satisfies $\hat s_t(x)\ge \beta_{\rm end}$. Since
$s_t(x)\ge \hat s_t(x)-\delta_s$ and $\psi_{\rm sp}$ is increasing, the claim
follows.
\end{proof}

\subsection{Exact Conditional-Risk Modulus}

\begin{definition}[Exact pointwise path defects]
\label{def:exact-pointwise-path-defects}
For each $x\in\mc{X}$, define
\begin{align}
  \delta_u^{\rm path}(x)
  &:=
  \sup_{t\in[0,1]}\abs{u_t(x)-\hat u_t(x)},
  \label{eq:delta-u-path}
  \\
  \delta_s^{\rm path}(x)
  &:=
  \sup_{t\in[0,1]}\abs{s_t(x)-\hat s_t(x)}.
  \label{eq:delta-s-path}
\end{align}
\end{definition}

\begin{proposition}[Pointwise exact-path raw-space barrier]
\label{prop:xdependent-exact-path-rawspace}
Assume that, for every $x\in\mc{X}$, the endpoints satisfy
\[
  \abs{u_{\Theta_A}(x)},\ \abs{u_{\Theta_B}(x)}\le M_{\rm end}(x),
  \qquad
  -\eta_{2,\Theta_A}(x),\ -\eta_{2,\Theta_B}(x)\ge \lambda_{\rm end}(x),
  \qquad
  \lambda_{\rm end}(x)>\lambda_{\min}.
\]
Under Definition~\ref{def:exact-pointwise-path-defects}, set
\begin{align*}
  \beta_{\rm end}(x)
  &:=
  \psi_{\rm sp}^{-1}(\lambda_{\rm end}(x)),
  \\
  M_{\rm path}(x)
  &:=
  M_{\rm end}(x)+\delta_u^{\rm path}(x),
  \\
  \lambda_{\rm path}(x)
  &:=
  \psi_{\rm sp}(\beta_{\rm end}(x)-\delta_s^{\rm path}(x)).
\end{align*}
Let
\[
  m_\ast(x):=\E[Y\mid X=x],
  \qquad
  m_{2,\ast}(x):=\E[Y^2\mid X=x].
\]
Define
\begin{align}
  C_u^{\rm cond}(x)
  &:=
  \frac{M_{\rm path}(x)}{2\lambda_{\rm path}(x)}+\abs{m_\ast(x)},
  \label{eq:Cu-cond}
  \\
  C_s^{\rm cond}(x)
  &:=
  \sup_{r\ge \lambda_{\rm path}(x)}
  \Bigl(1-e^{-(r-\lambda_{\min})}\Bigr)
  \max\Bigl\{
    \abs{\frac{1}{2r}-m_{2,\ast}(x)},
    \abs{\frac{M_{\rm path}(x)^2}{4r^2}+\frac{1}{2r}-m_{2,\ast}(x)}
  \Bigr\},
  \label{eq:Cs-cond}
  \\
  C_2^{\rm cond}(x)
  &:=
  \sup_{r\ge \lambda_{\rm path}(x)}
  \max\Bigl\{
    \abs{\frac{1}{2r}-m_{2,\ast}(x)},
    \abs{\frac{M_{\rm path}(x)^2}{4r^2}+\frac{1}{2r}-m_{2,\ast}(x)}
  \Bigr\}.
  \label{eq:C2-cond}
\end{align}
Then
\begin{equation}
  \cB(\Theta_A,\Theta_B;\varpi)
  \le
  \E\Bigl[
    C_u^{\rm cond}(X)\,\delta_u^{\rm path}(X)
    +
    C_s^{\rm cond}(X)\,\delta_s^{\rm path}(X)
    +
    C_2^{\rm cond}(X)\,\frac{(s_{\Theta_A}(X)-s_{\Theta_B}(X))^2}{32}
  \Bigr].
  \label{eq:xdependent-exact-path-rawspace}
\end{equation}
\end{proposition}

\begin{proof}
Fix $x$. The endpoint strip and the path defects imply
\[
  \abs{u_t(x)}\le M_{\rm path}(x),
  \qquad
  -\eta_{2,\Theta_t}(x)\ge \lambda_{\rm path}(x).
\]
Condition on $X=x$ and apply the mean value theorem directly to the raw-space
conditional risk
\[
  r_x(u,s):=\E[\tilde\ell(u,s;Y)\mid X=x].
\]
The final term is the pointwise softplus Jensen contribution.
\end{proof}

\begin{proposition}[Timewise exact conditional-risk modulus]
\label{prop:timewise-exact-condrisk-modulus}
Assume the hypotheses of Proposition~\ref{prop:xdependent-exact-path-rawspace}.
For each $t\in[0,1]$ and $x\in\mc{X}$, let
\[
  \delta_{u,t}(x):=\abs{u_t(x)-\hat u_t(x)},
  \qquad
  \delta_{s,t}(x):=\abs{s_t(x)-\hat s_t(x)},
  \qquad
  \hat r_t(x):=\psi_{\rm sp}(\hat s_t(x)).
\]
Define
\[
  I_t(x)
  :=
  \Bigl[
    \psi_{\rm sp}\bigl(\hat s_t(x)-\delta_{s,t}(x)\bigr),
    \psi_{\rm sp}\bigl(\hat s_t(x)+\delta_{s,t}(x)\bigr)
  \Bigr].
\]
Write
\begin{align}
  \Omega_{u,t}(x)
  &:=
  \sup_{r\in I_t(x)}
  \left\{
    \delta_{u,t}(x)\abs{\frac{\hat u_t(x)}{2r}-m_\ast(x)}
    +
    \frac{\delta_{u,t}(x)^2}{4r}
  \right\},
  \label{eq:omega-u}
  \\
  \Omega_{s,t}(x)
  &:=
  \sup_{r\in I_t(x)}
  \max\left\{
    0,\,
    m_{2,\ast}(x)\bigl(r-\hat r_t(x)\bigr)
    +
    \frac{\hat u_t(x)^2}{4}
    \left(
      \frac{1}{r}-\frac{1}{\hat r_t(x)}
    \right)
    +
    \frac{1}{2}\log\frac{\hat r_t(x)}{r}
  \right\},
  \label{eq:omega-s}
  \\
  J_t(x)
  &:=
  \max\left\{
    0,\,
    r_x\bigl(\hat u_t(x),\hat s_t(x)\bigr)
    -
    (1-t)\,r_x\bigl(u_{\Theta_A}(x),s_{\Theta_A}(x)\bigr)
    -
    t\,r_x\bigl(u_{\Theta_B}(x),s_{\Theta_B}(x)\bigr)
  \right\}.
  \label{eq:Jt}
\end{align}
Then
\begin{equation}
  \cB(\Theta_A,\Theta_B;\varpi)
  \le
  \sup_{t\in[0,1]}
  \E\bigl[
    \Omega_{u,t}(X)+\Omega_{s,t}(X)+J_t(X)
  \bigr].
  \label{eq:timewise-exact-condrisk-modulus}
\end{equation}
\end{proposition}

\begin{proof}
Condition on $X=x$ and split
\begin{align*}
  r_x\bigl(u_t(x),s_t(x)\bigr)-r_x\bigl(\hat u_t(x),\hat s_t(x)\bigr)
  &=
  \Bigl[
    r_x\bigl(u_t(x),s_t(x)\bigr)-r_x\bigl(\hat u_t(x),s_t(x)\bigr)
  \Bigr]
  \\
  &\quad+
  \Bigl[
    r_x\bigl(\hat u_t(x),s_t(x)\bigr)-r_x\bigl(\hat u_t(x),\hat s_t(x)\bigr)
  \Bigr].
\end{align*}
The first bracket is a quadratic increment in $u$ at fixed precision
$r=\psi_{\rm sp}(s_t(x))$, hence bounded by \eqref{eq:omega-u}. The second
bracket is the exact scale increment of the conditional risk at fixed
$\hat u_t(x)$, hence bounded by \eqref{eq:omega-s}. The endpoint interpolation
gap contributes \eqref{eq:Jt}. Taking $\E$ and $\sup_t$ yields the result.
\end{proof}

\begin{corollary}[Envelope exact conditional-risk modulus]
\label{cor:envelope-exact-condrisk-modulus}
Under the hypotheses of
Proposition~\ref{prop:xdependent-exact-path-rawspace}, define
\[
  I_t^{\rm env}(x)
  :=
  \Bigl[
    \psi_{\rm sp}\bigl(\hat s_t(x)-\delta_s^{\rm path}(x)\bigr),
    \psi_{\rm sp}\bigl(\hat s_t(x)+\delta_s^{\rm path}(x)\bigr)
  \Bigr].
\]
Let $\Omega^{\rm env}_{u,t}(x)$, $\Omega^{\rm env}_{s,t}(x)$ be obtained from
\eqref{eq:omega-u}--\eqref{eq:omega-s} by replacing
\[
  \delta_{u,t}(x),\ \delta_{s,t}(x),\ I_t(x)
  \quad\text{with}\quad
  \delta_u^{\rm path}(x),\ \delta_s^{\rm path}(x),\ I_t^{\rm env}(x).
\]
Then
\begin{equation}
  \cB(\Theta_A,\Theta_B;\varpi)
  \le
  \sup_{t\in[0,1]}
  \E\bigl[
    \Omega^{\rm env}_{u,t}(X)
    +
    \Omega^{\rm env}_{s,t}(X)
    +
    J_t(X)
  \bigr].
  \label{eq:envelope-exact-condrisk-modulus}
\end{equation}
\end{corollary}

\begin{proof}
This is immediate from
\[
  \delta_{u,t}(x)\le\delta_u^{\rm path}(x),
  \qquad
  \delta_{s,t}(x)\le\delta_s^{\rm path}(x),
\]
and the interval inclusion
\[
  I_t(x)\subseteq I_t^{\rm env}(x),
\]
which follows from the monotonicity of $\psi_{\rm sp}$.
\end{proof}

\begin{remark}[Final sharpened route]
The appendix-level route is now direct:
\begin{align*}
  &\text{exact projected transfer}
  \;\to\;
  \text{endpoint conditional-risk strip}
  \\
  &\to\;
  \text{raw-scale propagation}
  \;\to\;
  \text{exact conditional-risk modulus}.
\end{align*}
Proposition~\ref{prop:exact-sigmoid-projected-transfer} supplies the transfer
terms. Proposition~\ref{prop:conditional-risk-strip} and
Proposition~\ref{prop:raw-scale-strip-propagation} produce the propagated
variance strip. Finally,
Corollary~\ref{cor:envelope-exact-condrisk-modulus} or
Proposition~\ref{prop:timewise-exact-condrisk-modulus} closes the barrier
bound.
\end{remark}

\subsection{Quantitative Instantiation}

The final gap is governed by how much of the loss increment is still treated by
a worst-case Lipschitz reduction. The exact conditional-risk modulus removes
that last source of slack.

\begin{remark}[Quantitative status on the matched softplus run]
\label{rem:quantitative-status}
On the matched softplus run in
\texttt{experiments/results\_softplus\_sharp3/natural}, the relevant
theorem-consistent instantiated certificates are:
\[
  \cB \lesssim 1.24\times 10^1
  \quad\text{(exact projected transfer + raw-space barrier),}
\]
\[
  \cB \lesssim 1.64
  \quad\text{(endpoint exact projected raw-space barrier),}
\]
\[
  \cB \lesssim 1.53\times 10^{-2}
  \quad\text{(envelope exact conditional-risk modulus),}
\]
\[
  \cB \lesssim 1.43\times 10^{-2}
  \quad\text{(timewise exact conditional-risk modulus).}
\]
The observed matched barrier is
\[
  \cB_{\rm obs}\approx 1.01\times 10^{-2}.
\]
Thus the honest theorem-consistent instantiated-certificate gap is now only
about a factor of $1.4$--$1.5$.
\end{remark}

\begin{remark}[What the final decomposition says]
The route totals in Remark~\ref{rem:quantitative-status} are
$\sup_t\E[\Omega_{u,t}+\Omega_{s,t}+J_t]$ or its envelope analogue. By
contrast, the components displayed below are per-term maxima
$\sup_t\E[\Omega_{u,t}]$, $\sup_t\E[\Omega_{s,t}]$, and $\sup_t\E[J_t]$, so
they need not sum exactly to the reported route totals.

For the timewise exact-modulus route, the per-term maxima are
\[
  \text{$u$-modulus}\approx 9.93\times 10^{-3},
  \qquad
  \text{$s$-modulus}\approx 4.38\times 10^{-3},
  \qquad
  J_t\approx 6.2\times 10^{-8}.
\]
For the envelope route, the analogous per-term maxima are
\[
  \text{$u$-envelope}\approx 1.05\times 10^{-2},
  \qquad
  \text{$s$-envelope}\approx 4.94\times 10^{-3},
  \qquad
  J_t\approx 6.2\times 10^{-8}.
\]
So at this point neither Jensen nor transfer dominates by orders of magnitude.
The remaining gap is mild constant-factor slack in the exact conditional-risk
modulus itself.
\end{remark}

\begin{thebibliography}{99}

\bibitem{Ferbach2024}
Damien Ferbach, Baptiste Goujaud, Gauthier Gidel, and Aymeric Dieuleveut.
\newblock Proving linear mode connectivity of neural networks via optimal transport.
\newblock In \emph{Proceedings of The 27th International Conference on Artificial Intelligence and Statistics (AISTATS)}, volume 238 of \emph{Proceedings of Machine Learning Research}, pages 3853--3861, 2024.

\bibitem{Frankle2020}
Jonathan Frankle, Gintare Karolina Dziugaite, Daniel Roy, and Michael Carbin.
\newblock Linear mode connectivity and the lottery ticket hypothesis.
\newblock In \emph{Proceedings of the 37th International Conference on Machine Learning (ICML)}, pages 3259--3269, 2020.

\bibitem{Garipov2018}
Timur Garipov, Pavel Izmailov, Dmitrii Podoprikhin, Dmitry Vetrov, and Andrew Gordon Wilson.
\newblock Loss surfaces, mode connectivity, and fast ensembling of {DNN}s.
\newblock In \emph{Advances in Neural Information Processing Systems 31 (NeurIPS)}, pages 8789--8798, 2018.

\bibitem{MeiMisiakiewiczMontanari2019}
Song Mei, Theodor Misiakiewicz, and Andrea Montanari.
\newblock Mean-field theory of two-layer neural networks: dimension-free bounds and kernel limit.
\newblock In \emph{Proceedings of the 32nd Conference on Learning Theory (COLT)}, volume 99 of \emph{Proceedings of Machine Learning Research}, pages 2388--2464, 2019.

\bibitem{WainwrightJordan2008}
Martin J. Wainwright and Michael I. Jordan.
\newblock Graphical models, exponential families, and variational inference.
\newblock \emph{Foundations and Trends in Machine Learning}, 1(1--2):1--305, 2008.

\end{thebibliography}

\end{document}
